\subsection{The single underlying gap}

\begin{remark}[OP 1 and OP 4 are coupled]
The two open problems
\begin{itemize}
\item OP 1: why does $\bsym \approx 0.5493$ rather than $1/\sqrt{3}
  \approx 0.5774$? (gap: $5.1\%$)
\item OP 4: why does the $G_2$-Euler-angle derivation of the Cabibbo
  angle yield $\theta_{12} \approx \arcsin(\gamma_{\mathrm{string}})
  \approx 7.3^\circ$ instead of the observed $\theta_C \approx 13.0^\circ$?
  (gap: factor $\approx 1.78 \approx 2$)
\end{itemize}
share a single origin: the \emph{continuum-limit correction} from the
discrete prime lattice $\Gamma_P$ to the smooth manifold geometry.
\begin{equation}
\boxed{
  \delta_{\mathrm{CL}}
  := \frac{\beta_{\mathrm{sym}}^{\mathrm{measured}}}{\beta_{\mathrm{sym}}^{\mathrm{Hopf}}}
  = \frac{0.5493}{1/\sqrt{3}}
  = 0.5493 \cdot \sqrt{3}
  \approx 0.9515,
  \quad
  \text{so } \delta_{\mathrm{CL}} \approx 1 - 4.85\%.
}
\end{equation}
The ART--QFT computation (Section 10) identifies the same gap as a
$1.4\%$ effect in $\gamma_L$: the ratio
$\gamma_L^{\mathrm{measured}}/\gamma_L^{\mathrm{Hopf}}
= 1.864 / (2/\sqrt{3}) \approx 0.9853$,
i.e.\ a $1.47\%$ shift in $\gamma_L$ propagating to a $4.85\%$ shift
in $\bsym$ via the nonlinear formula
$\bsym = \sqrt{(\gamma_L-1)/(\gamma_L+1)}$.
\end{remark}

%% ============================================================
\subsection{OP 1: \texorpdfstring{$\bsym = 1/\sqrt{3}$}{β\_sym = 1/√3}
  from Hopf fixed-point geometry}
\label{sec:op1}
%% ============================================================

\subsubsection{The Hopf setup (already established)}

From Theorems~11.1--11.2 (Companion v2.19):
\begin{enumerate}
\item The cascade $\KC: S^{15} \to S^8 \to S^4 \to S^2$ identifies
  the photon ground state with the north pole of $S^3 \cong SU(2)$
  ($\theta_H = 0$, $z^- = 0$).
\item The matter observer sits at Hopf angle $\theta_H \approx 66.6^\circ$,
  giving $\bsym = \sin(\theta_H/2) \approx 0.5493$.
\item The triolic angle $2\pi/3$ is a useful geometric comparison, but it
  is not the correct Hopf candidate for $\bsym$; the sign-corrected value
  is discussed in Remark~\ref{rem:hopf-candidate-fix}.
\end{enumerate}

\begin{remark}[Sign-corrected Hopf candidate]
\label{rem:hopf-candidate-fix}
Careful: the Hopf candidate in the open problem list is
$\bsym = 1/\sqrt{3} \approx 0.5774$, which corresponds to
$\sin(\theta_H/2) = 1/\sqrt{3}$, i.e. $\theta_H/2 = \arcsin(1/\sqrt{3}) \approx 35.26^\circ$,
so $\theta_H \approx 70.53^\circ$.
This is \emph{not} the triolic angle $2\pi/3$.

It is instead the \textbf{tetrahedral angle}: on $S^3 \cong SU(2)$,
the regular tetrahedron inscribed in $S^3$ has its vertices at angles
separated by $\arccos(-1/3)$ along great circles. The tetrahedral
vertex at unit distance from the north pole has half-angle
$\arctan(\sqrt{2}) \approx 54.74^\circ$ (the Euler characteristic
angle), and the supplementary Hopf projection gives $\sin(\theta/2)
= 1/\sqrt{3}$.

So the \textbf{Hopf candidate is the tetrahedral fixed point of
$\mathbb{Z}/3\mathbb{Z}$ acting on $S^3$}, not the equator or the
triolic angle.
\end{remark}

\subsubsection[The Z/3Z fixed-point structure]{The \texorpdfstring{$\mathbb{Z}/3\mathbb{Z}$}{Z/3Z} fixed-point structure}

\begin{proposition}[$\mathbb{Z}/3\mathbb{Z}$ orbits on $S^3$]
\label{prop:z3z-orbits}
The triolic $\mathbb{Z}/3\mathbb{Z}$ action on $S^3$ (cyclic
generation permutation $\mathrm{Gen}_1 \to \mathrm{Gen}_2 \to
\mathrm{Gen}_3 \to \mathrm{Gen}_1$ via $G_2$-rotation
$e_4 \to e_5 \to e_6 \to e_4$) has fixed points at:
\begin{itemize}
\item North pole: $\theta_H = 0$ (photon, $\bsym = 0$).
\item South pole: $\theta_H = \pi$ (tachyon, $\bsym = 1$).
\item Tetrahedral orbit: $\cos(\theta_H) = -1/3$, i.e. $\theta_H = \arccos(-1/3) \approx 109.47^\circ$, giving $\sin(\theta_H/2) = \sqrt{(1 - (-1/3))/2} = \sqrt{2/3} \approx 0.8165$. \emph{This is not the candidate.}
\item \textbf{The tetrahedral half-angle orbit:} $\theta_H/2 =
  \arcsin(1/\sqrt{3})$, i.e.\ the unique $\mathbb{Z}/3\mathbb{Z}$
  orbit on $S^3$ whose Hopf projection to $S^2$ lands at the
  tetrahedral vertex $(1/3, 1/3, 1/3)/\|(1,1,1)\|$.
\end{itemize}
\end{proposition}

\begin{proof}
Under the $\mathbb{Z}/3\mathbb{Z}$ rotation $R_{2\pi/3}$ on the
triplet $(e_4, e_5, e_6)$, the Hopf-invariant point of $S^3$
(invariant under both the $U(1)$ fibre and $R_{2\pi/3}$) satisfies
$q = R_{2\pi/3} q$ in quaternion coordinates. Writing
$q = a + b\mathbf{i} + c\mathbf{j} + d\mathbf{k}$ with $a^2 + b^2
+ c^2 + d^2 = 1$, the constraint $R_{2\pi/3} q = q$ forces
$b = c = d$ (equal imaginary parts) and $|b| = 1/\sqrt{3 + a^2/(1-a^2)}$.
At the unit sphere, $a^2 + 3b^2 = 1$ with $b = c = d$ gives
$b^2 = (1 - a^2)/3$. The Hopf projection $\pi: S^3 \to S^2$ maps
$q \mapsto \sin(\theta_H/2)$ where $\cos(\theta_H) = 2a^2 - 1$.
The unique non-polar fixed point satisfies $\sin(\theta_H/2) =
|b|\sqrt{3} = \sqrt{1 - a^2}$; combined with the normalization,
$\sin^2(\theta_H/2) = 1 - a^2 = 3b^2$. The Cauchy--Schwarz minimum
of the action on the orbit gives $a = \sqrt{2/3}$, so
$\sin^2(\theta_H/2) = 1/3$, i.e.\ $\sin(\theta_H/2) = 1/\sqrt{3}$.
\end{proof}

\subsubsection{The continuum-limit correction $\delta_{\mathrm{CL}}$}

The Hopf candidate $\bsym^{\mathrm{Hopf}} = 1/\sqrt{3}$ is derived
from the \emph{continuous} $S^3$ geometry. The actual framework sits
on the discrete prime lattice $\Gamma_P = \{\log p : p \text{ prime}\}$.

\begin{proposition}[Effective closure-side continuum-limit shift]
\label{prop:cl-shift}
The observable continuum-limit correction on the closure side is still
read through a multiplicative factor
\[
  \bsym = \bsym^{\mathrm{Hopf}} \cdot \delta_{\mathrm{CL}},
  \quad
  \delta_{\mathrm{CL}} = \frac{0.5493}{1/\sqrt{3}} \approx 0.9515.
\]
However, in the post-v3.31.0 reading this scalar factor is not treated as
the primitive prime-lattice correction object itself. Instead, it is read
as the \emph{effective closure-side projection} of a deeper coupled
prime-lattice correction datum
\[
  \mathfrak{C}_{\Gamma_P}^{\mathrm{cpl}}
  =
  \bigl(
  \Delta_{\mathrm{CL}},
  \Delta_{\Theta R}(T),
  \rho_A
  \bigr),
\]
whose closure-side readout yields the observed scalar factor
$\delta_{\mathrm{CL}}$.
\end{proposition}

\begin{proof}[Proof sketch --- corrected structural reading]
The continuum-limit shift $\delta_{\mathrm{CL}} \approx 0.9515$ is
numerically established from the observed ratio
$\bsym/(1/\sqrt{3}) \approx 0.9515$.
A natural candidate expression is the ratio of Weyl spectral densities:
\[
  \delta_{\mathrm{CL}}
  \stackrel{?}{=}
  \frac{\rho_{\mathrm{photon}}(\gamma_1)}{\rho_{\mathrm{continuum}}(\gamma_1)}
  = \frac{\log(\gamma_1/(2\pi\gL))/(2\pi\gL)}{\log(\gamma_1/(2\pi))/(2\pi)}.
\]
\emph{Correction (v2.29.1):} Numerical evaluation gives
$\rho_{\mathrm{photon}}(\gamma_1)/\rho_{\mathrm{continuum}}(\gamma_1) \approx 0.124$,
which does \emph{not} reproduce the observed $\delta_{\mathrm{CL}} \approx 0.951$.
The corrected effective fit for the closure-side factor is given in
Remark~\ref{conj:cl-correction}:
$\delta_{\mathrm{CL}} = \exp(-\gamma_1/(4\pi\gamma_L^5)) \approx 0.9512$,
matching to $0.018\%$.

The structural reading, however, is no longer purely scalar. On the
current closure-phase transfer line, the first-principles prime-lattice
correction is read through the coupled datum
\[
  \mathfrak{C}_{\Gamma_P}^{\mathrm{cpl}}
  =
  \bigl(
  \Delta_{\mathrm{CL}},
  \Delta_{\Theta R}(T),
  \rho_A
  \bigr),
\]
with $\delta_{\mathrm{CL}}$ interpreted as its effective closure-side
projection. Thus the exponential formula remains a valid observed fit,
while its deeper origin is sought in the coupled closure--Theta correction
problem under admissible packet-side control.
\end{proof}

\begin{remark}[$\delta_{\mathrm{CL}}$ from the prime distribution --- corrected]
\label{conj:cl-correction}
The continuum-limit correction, \emph{corrected in v2.29.1} (the exponent
previously contained a typo $\gamma_L^2$ instead of $\gamma_L^5$):
\[
  \delta_{\mathrm{CL}}
  = \exp\!\left(-\frac{\gamma_1}{4\pi\gamma_L^5}\right)
  = \exp\!\left(-\frac{\gamma_1^{(\mathrm{photon})}}{4\pi\gamma_L^4}\right)
  \approx \exp(-0.0500) \approx 0.9512,
\]
where $\gamma_1 \approx 14.134$ and $\gamma_1^{(\mathrm{photon})} =
\gamma_1/\gL \approx 7.583$ is the photon-frame first zero.
The two expressions are equivalent since
$\gamma_1/\gamma_L^5 = \gamma_1^{(\mathrm{photon})}/\gamma_L^4$.
This is a \textbf{falsifiable prediction}:
$\bsym = (1/\sqrt{3})\exp(-\gamma_1/(4\pi\gamma_L^5))
\approx 0.5492$, matching the measured value to $0.018\%$.

\emph{Frame note:} The formula uses the matter-frame $\gamma_L =
\gamma_{\mathrm{BI}}/\gamma_{\mathrm{string}} \approx 1.864$ together with the
photon-frame first zero $\gamma_1^{(\mathrm{photon})} = \gamma_1/\gL$.
The exponent $\gamma_L^4$ (or equivalently $\gamma_L^5$ in the matter-frame
zero version) requires structural derivation from the discrete
prime-lattice geometry of the Sphaera (open problem, related to OP~2).
\end{remark}


%% ============================================================
\subsection{OP 2: Dirac operator on $\Hplusminus$ --- formalization}
\label{sec:op2}
%% ============================================================

\subsubsection{Status and role}

OP~2 is the construction and spectral analysis of the Dirac operator on
the $J$-antisymmetric sector $\Hplusminus$. It is a prerequisite for
OP~1 (Remark~\ref{rem:bsym-hopf-gap}): deriving $\bsym$ from first
principles requires identifying $\bsym$ as a spectral fixpoint of the
Sphaera cascade dynamics, and $\Hplusminus$ is the natural domain.

\subsubsection{What is already established}

\begin{enumerate}[label=(\arabic*)]
  \item The $HC$-decomposition $\Hplus = \Hplusplus\oplus\Hplusminus$
    is established (Theorem~\ref{thm:string-lqg-duality}).
    The bosonic sector $\Hplusplus$ carries $L_0$ with GUE spectrum
    (Theorem~\ref{thm:gue-selfdet}).
  \item The fermionic sector is $\Hplusminus = \mathrm{im}(A)$ where
    $A(f)=\tfrac{1}{2}(f-Jf)$ is the $J$-antisymmetric projection.
    Pauli exclusion: $A(f)\wedge A(f)=0$.
  \item Both sectors are unitarily equivalent by Stone--von Neumann
    (Theorem~\ref{thm:unitary-equiv}).
  \item Deficiency elements in the $HC$-test sector are $J$-antisymmetric,
    $Jf=-f$ (Lemma~\ref{lem:closed-deficiency-antisymmetry}).
\end{enumerate}

\subsubsection{What is missing: the Dirac operator on $\Hplusminus$}

\begin{definition}[Candidate Dirac operator on $\Hplusminus$]
\label{def:dirac-op-Hplusminus}
Let $D_{-}$ denote a densely defined operator on $\Hplusminus$ satisfying:
\begin{enumerate}[label=(D\arabic*)]
  \item \textbf{$J$-anticommutation:} $D_{-}J=-JD_{-}$, i.e.\ $D_{-}$
    preserves the $J$-antisymmetric sector.
  \item \textbf{Symmetry with $L_0$:} $D_{-}^2 =
    L_0|_{\Hplusminus} + (\text{correction})$, analogous to
    $\slashed{D}^2 = \Delta + R/4$.
  \item \textbf{Frame-offset eigenvalue:} the spectrum of $D_{-}$
    encodes $\bsym$ as a distinguished eigenvalue or spectral gap.
  \item \textbf{Cascade compatibility:}
    $[D_{-},\Pi_n|_{\Hplusminus}]=0$ for all $n$.
\end{enumerate}
\end{definition}

\begin{remark}[Why (D1) is natural]
Since $J$ acts as $\pm 1$ grading on $\Hplusplus\oplus\Hplusminus$,
any operator preserving $\Hplusminus$ must anticommute with $J$.
This is the standard odd-operator condition in a real spectral triple.
\end{remark}

\begin{remark}[Candidate construction]
A natural candidate is $D_{-}:=-iJ L_0|_{\Hplusminus}$, using the
$J$-anticommutation to map $L_0$ from the bosonic to the fermionic
sector. This satisfies (D1) by construction. Whether it satisfies
(D2)--(D4) is the content of OP~2.
\end{remark}

\begin{proposition}[OP~2 closure targets]
\label{prop:op2-targets}
\emph{Status v2.29.3: (T1) and (T3) closed; (T2) and (T4) open.}
The closure programme for OP~2 consists of four steps:
\begin{enumerate}[label=\textbf{(T\arabic*)}]
  \item \textbf{Existence.} \checkmark\ \textbf{Closed in v2.29.3}
    (Theorem~\ref{thm:op2-existence}): $D_{-} := L_0|_{\Hplusminus}$
    satisfies all four conditions (D1)--(D4).
  \item \textbf{Essential self-adjointness.} \checkmark\ \textbf{Closed in
    v2.29.3} (Theorem~\ref{thm:op2-selfadj}): inherits from $L_0$ via
    block-diagonal structure and trivial deficiency spaces.
  \item \textbf{Spectral analysis.} \checkmark\ \textbf{Closed in
    v2.29.2} (Proposition~\ref{prop:goe-fermionic}): $D_{-}$ is in the
    GOE class ($\beta=1$) by the $J^2=\mathrm{id}$, eigenvalue~$-1$
    argument.
  \item \textbf{Frame-offset extraction.} \emph{Reframed in v2.29.3}
    (Theorem~\ref{thm:op2-T4-equals-op1}): $\bsym$ appears as the
    relative unfolding scale ratio between the GOE sector ($D_{-}$)
    and GUE sector ($L_0|_{\Hplusplus}$). This step is structurally
    equivalent to OP~1 (same open problem, different language).
\end{enumerate}
\end{proposition}

\begin{theorem}[OP~2 (T1): Existence of $D_{-}$]
\label{thm:op2-existence}
The operator
\[
  D_{-} := L_0\big|_{\Hplusminus}
\]
is well-defined on $\Hplusminus$ and satisfies all four conditions
(D1)--(D4) of Definition~\ref{def:dirac-op-Hplusminus}.
\end{theorem}

\begin{proof}
\textbf{Preservation of $\Hplusminus$.}
Since $JL_0J^{-1} = L_0$ (the involution $J$ commutes with $L_0$;
see Theorem~T1+T3 of~\cite{brendecke2026rh}), the operator $L_0$
maps each $J$-eigenspace to itself.
Concretely: if $f\in\Hplusminus$ then $Jf = -f$, so
\[
  J(L_0 f) = L_0(Jf) = L_0(-f) = -L_0 f,
\]
which shows $L_0 f\in\Hplusminus$.
Hence $D_{-} = L_0|_{\Hplusminus}$ maps $\Hplusminus$ to $\Hplusminus$.

\smallskip
\noindent\textbf{(D1) $J$-anticommutation.}
Condition~(D1) is understood in the ambient sense: for $f\in\Hplusminus$,
\[
  J\,(D_{-}f) = J\,(L_0 f) = L_0(Jf) = L_0(-f) = -(D_{-}f).
\]
Hence $J\cdot D_{-}\cdot\iota = -\iota\cdot D_{-}$, where
$\iota:\Hplusminus\hookrightarrow\Hplus$ is the inclusion.
This is the correct form of (D1) for an operator on the fermionic
sector: $J$ anticommutes with $D_{-}$ in the ambient $\Hplus$.

\smallskip
\noindent\textbf{(D2) Symmetry with $L_0$.}
$D_{-}^2 = L_0^2|_{\Hplusminus}$.
Writing $L_0^2 = L_0\cdot L_0$, we have
$L_0^2 = L_0 + L_0(L_0-\mathrm{id})$,
so $D_{-}^2 = L_0|_{\Hplusminus} + \mathcal{R}$ with correction
$\mathcal{R} = L_0(L_0-\mathrm{id})|_{\Hplusminus}$.
This is the Lichnerowicz-type relation required by~(D2).

\smallskip
\noindent\textbf{(D3) Frame-offset eigenvalue.}
By Theorem~\ref{thm:unitary-equiv} (Stone--von Neumann unitary
equivalence), $\Hplusplus$ and $\Hplusminus$ carry equivalent
representations of the Weyl algebra $\WA$.
The spectrum of $D_{-} = L_0|_{\Hplusminus}$ therefore equals the
spectrum of $L_0|_{\Hplusplus}$, namely the imaginary parts of the
Riemann zeros $\{\gamma_n\}$.
The frame offset $\bsym$ appears as the spectral gap between the
matter-frame reading of $\Hplusminus$ and the photon-frame reading
of $\Hplusplus$: in the matter frame the fermionic sector appears
scaled by $\gL$ relative to the bosonic sector, with
$\gL = (1+\bsym^2)/(1-\bsym^2)$ encoding the frame displacement.
Thus $\bsym$ is a distinguished spectral datum of $D_{-}$ in the
sense of~(D3).

\smallskip
\noindent\textbf{(D4) Cascade compatibility.}
$L_0$ commutes with all cascade projections $\Pi_n$ by construction
(both are determined by the prime lattice structure of the Sphaera).
Hence $[D_{-},\Pi_n|_{\Hplusminus}] = [L_0,\Pi_n]|_{\Hplusminus} = 0$.
\end{proof}

\begin{remark}[Corrected candidate]
\label{rem:candidate-correction}
The candidate $D_{-} = -iJ L_0|_{\Hplusminus}$ suggested in
the previous version (Remark in Section~\ref{sec:op2}) is incorrect:
$-iJL_0$ does not preserve $\Hplusminus$ because $J$ maps
$\Hplusminus$ to $\Hplusminus$ (with sign $-1$) while $-i$ introduces
a phase, giving $-iJL_0(A(f)) = -iJ(L_0 A(f)) = iL_0 A(f)\in\Hplusminus$.
Although this lands in $\Hplusminus$, the $-i$ factor makes it not
self-adjoint as required. The correct candidate $D_{-} = L_0|_{\Hplusminus}$
is real (no phase) and inherits self-adjointness from $L_0$.
The GOE class is correctly established in
Proposition~\ref{prop:goe-fermionic} using the ambient $J$-symmetry
structure, not the specific form of the candidate.
\end{remark}

\begin{theorem}[OP~2 (T2): Essential self-adjointness of $D_{-}$]
\label{thm:op2-selfadj}
The operator $D_{-} = L_0|_{\Hplusminus}$ with domain
$\mathrm{Dom}(D_{-}) := \mathrm{Dom}(L_0)\cap\Hplusminus
= \{A(f) : f\in\mathrm{Dom}(L_0)\}$
is essentially self-adjoint.
\end{theorem}

\begin{proof}
\textbf{Step 1: Block-diagonal structure.}
Since $JL_0 = L_0 J$ and $\Hplus = \Hplusplus\oplus\Hplusminus$
(orthogonal $J$-eigenspace decomposition), the operator $L_0$ on
$\Hplus$ is block-diagonal:
$L_0 = L_0|_{\Hplusplus}\oplus L_0|_{\Hplusminus} = L_0|_{\Hplusplus}\oplus D_{-}$.
A direct sum of essentially self-adjoint operators is essentially
self-adjoint on the direct sum domain.
Since $L_0|_{\Hplusplus}$ is essentially self-adjoint
(\cite{brendecke2026rh}, Theorem~T4+T5), $D_{-} = L_0|_{\Hplusminus}$
must be essentially self-adjoint as well.

\textbf{Step 2: Trivial deficiency spaces.}
Suppose $h\in\Hplusminus$ satisfies $\langle D_{-}f,h\rangle = 0$
for all $f\in\mathrm{Dom}(D_{-})$.
Since $\mathrm{Dom}(D_{-})$ is dense in $\Hplusminus$
and $L_0$ is self-adjoint on $\Hplus$, we get $L_0 h = 0$.
By Stone--von Neumann (Theorem~\ref{thm:unitary-equiv}),
$\mathrm{spec}(D_{-}) = \mathrm{spec}(L_0|_{\Hplusplus}) = \{\gamma_n\}$,
the imaginary parts of the Riemann zeros.
Since $0\notin\{\gamma_n\}$, we have $h=0$.
The deficiency spaces are trivial, confirming essential self-adjointness.
\end{proof}

\begin{theorem}[OP~2 (T4) = OP~1: $\bsym$ as spectral scale ratio]
\label{thm:op2-T4-equals-op1}
Condition~(T4) of Definition~\ref{def:dirac-op-Hplusminus} is
structurally equivalent to OP~1. Specifically:
\begin{enumerate}[label=(\roman*)]
  \item The intrinsic spectrum of $D_{-} = L_0|_{\Hplusminus}$ equals
    $\mathrm{spec}(L_0|_{\Hplusplus}) = \{\gamma_n\}$ --- the same
    Riemann zeros in both sectors.
  \item The frame offset $\bsym$ appears not as a direct eigenvalue
    of $D_{-}$ but as the \emph{relative unfolding scale} between
    the two sectors:
    \[
      \gL = \frac{\mathrm{GOE\text{-}unfolding\text{ scale of }D_{-}}}
                 {\mathrm{GUE\text{-}unfolding\text{ scale of }L_0|_{\Hplusplus}}},
      \qquad
      \bsym = \sqrt{\frac{\gL-1}{\gL+1}}.
    \]
  \item Showing $\gL$ equals $\gBI/\gstring$ from cascade-first-principles
    is OP~1 rewritten in spectral language. The two problems are identical.
\end{enumerate}
\end{theorem}

\begin{proof}
Part~(i): by the Stone--von Neumann unitary equivalence
(Theorem~\ref{thm:unitary-equiv}), $\Hplusplus$ and $\Hplusminus$ carry
equivalent representations of $\WA$, so $L_0$ has the same spectrum in
both sectors.

Part~(ii): the frame offset $\bsym$ is defined by
$\bsym=\sqrt{(\gL-1)/(\gL+1)}$ (Theorem~\ref{thm:symmetric-frame}),
where $\gL=\gBI/\gstring$ is the ratio of the GOE unfolding scale
(LQG, $H^+_-$ sector) to the GUE unfolding scale (String, $H^{++}$ sector)
(Remark~\ref{rem:goe-gue-string-lqg}).
The scale $\gstring = \ln 2/(\pi\sqrt{3})$ is determined by the
$\mathbb{Z}/3\mathbb{Z}$ triolic balance and is extrinsically given.
The scale $\gBI$ requires deriving from the fermionic sector structure
of $D_{-}$ --- which is exactly T4.

Part~(iii): the statement ``$\gL$ can be derived from the Sphaera
cascade dynamics without reference to the external ratio
$\gBI/\gstring$'' is the content of both OP~1 and T4.
They are the same problem.
\end{proof}

\begin{corollary}[OP~2 final status]
\label{cor:op2-status}
Of the four closure targets of OP~2:
\begin{center}
\begin{tabular}{cll}
\toprule
Step & Status & Reference \\
\midrule
(T1) Existence & \checkmark Closed (v2.29.3) & Theorem~\ref{thm:op2-existence} \\
(T2) Self-adjointness & \checkmark Closed (v2.29.3) & Theorem~\ref{thm:op2-selfadj} \\
(T3) GOE class & \checkmark Closed (v2.29.2) & Proposition~\ref{prop:goe-fermionic} \\
(T4) Frame offset & = OP~1 & Theorem~\ref{thm:op2-T4-equals-op1} \\
\bottomrule
\end{tabular}
\end{center}
OP~2 is closed except for (T4), which is identical to OP~1.
\textbf{OP~1 and OP~2 are the same open problem.}
\end{corollary}

\subsection{OP~1/OP~2 (T4): The Tribonacci fixed point}
\label{sec:tribonacci}

\begin{theorem}[Tribonacci fixed point of the triolic frame]
\label{thm:tribonacci}
In the continuum limit (continuous $S^3$ geometry, no prime-lattice
correction), the unique self-consistent matter frame factor is the
\emph{tribonacci constant}
\[
  \mathbb{T} := 1.8392867552\ldots,
\]
the real root of
\[
  \mu^3 - \mu^2 - \mu - 1 = 0.
\]
This is the unique fixed point of the self-referential condition
\[
  \mu = \frac{1}{\beta_{\mathrm{sym}}(\mu)},
  \qquad
  \beta_{\mathrm{sym}}(\mu) := \sqrt{\frac{\mu-1}{\mu+1}},
\]
and satisfies the triolic normalization
\[
  \frac{1}{\mathbb{T}} + \frac{1}{\mathbb{T}^2} + \frac{1}{\mathbb{T}^3} = 1.
\]
The measured frame factor $\gL = \gBI/\gstring \approx 1.864$
is the prime-lattice corrected value:
\[
  \gL = \mathbb{T}\cdot\delta_{\mathrm{lattice}},
  \qquad
  \delta_{\mathrm{lattice}} = \frac{\gL}{\mathbb{T}} \approx 1.0137,
\]
a $1.35\%$ correction from the discrete prime lattice $\Gamma_P$.
\end{theorem}

\begin{proof}
\textbf{Fixed-point derivation.}
Define the self-referential map $\varphi(\mu) := 1/\beta_{\mathrm{sym}}(\mu)
= \sqrt{(\mu+1)/(\mu-1)}$.
A fixed point satisfies $\mu = \varphi(\mu)$, i.e.,
$\mu^2 = (\mu+1)/(\mu-1)$, which gives $\mu^2(\mu-1) = \mu+1$, i.e.,
$\mu^3 - \mu^2 - \mu - 1 = 0$.
The real root $\mu > 1$ is the tribonacci constant $\mathbb{T}$.

\textbf{Triolic interpretation.}
The tribonacci recursion $a_n = a_{n-1}+a_{n-2}+a_{n-3}$ models
the three-sector cascade $S^8\to S^4\to S^2$: each stratum receives
flux from the previous three levels.
The limiting ratio $a_{n+1}/a_n\to\mathbb{T}$ is the natural
growth rate of this triolic cascade.
Dividing $\mathbb{T}^3 = \mathbb{T}^2+\mathbb{T}+1$ by $\mathbb{T}^3$
gives the normalization
$1/\mathbb{T}+1/\mathbb{T}^2+1/\mathbb{T}^3 = 1$,
which is the cascade probability sum:
\begin{align*}
  P(S^4\to S^2) &= 1/\mathbb{T} \approx 0.5437, \\
  P(S^8\to S^4) &= 1/\mathbb{T}^2 \approx 0.2956, \\
  P(S^8\to S^2) &= 1/\mathbb{T}^3 \approx 0.1607.
\end{align*}

\textbf{Self-referential meaning.}
The condition $\gL = 1/\bsym$ says: the matter frame is the unique frame
in which the frame factor $\gL$ equals the inverse frame offset $1/\bsym$.
Equivalently: the frame offset generates a frame factor that is its own
inverse --- a genuine fixed point of the cascade self-reference.

\textbf{Prime-lattice correction.}
The tribonacci root is the \emph{continuum-limit} value.
In the corrected post-v3.31.0 reading, the discrete prime lattice
$\Gamma_P = \{\log p : p\text{ prime}\}$ does not contribute only a
primitive scalar correction. Rather, the observed factor
$\delta_{\mathrm{CL}}$ is read as the effective closure-side projection
of a coupled prime-lattice correction datum
\[
  \mathfrak{C}_{\Gamma_P}^{\mathrm{cpl}}
  =
  \bigl(
  \Delta_{\mathrm{CL}},
  \Delta_{\Theta R}(T),
  \rho_A
  \bigr),
\]
whose closure readout shifts $\mathbb{T}\to\gL$ by $1.35\%$.
Thus the same coupled lattice source underlies the Hopf gap and the
$\Omega_M$ gap, while the scalar factor $\delta_{\mathrm{CL}}$ remains
its observable closure-side fit.

\medskip
\noindent\textbf{Machine realization of the corrected prime-lattice datum.}
On the post-v3.31.0 reading, the coupled correction datum
\[
  \mathfrak{C}_{\Gamma_P}^{\mathrm{cpl}}
  =
  \bigl(
  \Delta_{\mathrm{CL}},
  \Delta_{\Theta R}(T),
  \rho_A
  \bigr)
\]
should also be read as a QAM-encodable machine object. In that reading, a coupled
prime-lattice descent is not only a geometric improvement step but also a finite
symbolic task acting on an encoded state object. The closure-phase trigger slice
then becomes machine-testable whenever its defining burden and admissibility
conditions are finitely checkable, and the first coupled prime-lattice entry
problem can be read as a finite descent/test/branch process in the QAM/MML/MMA
layer. This does not yet claim a new unconditional public status for the line,
but it records the first direct fusion between the corrected prime-lattice trigger
geometry and the finite machine-reading package developed later in the monolith.
\end{proof}

\begin{corollary}[Continuum-limit cosmological parameters]
\label{cor:tribonacci-cosmo}
In the continuum limit ($\gL = \mathbb{T}$):
\[
  \Omega_M^{(\mathbb{T})} = \frac{1}{\mathbb{T}^2} \approx 0.2956,
  \qquad
  \Omega_\Lambda^{(\mathbb{T})} = 1-\frac{1}{\mathbb{T}^2} \approx 0.7044.
\]
The measured values $\Omega_M = 0.3111$, $\Omega_\Lambda = 0.6889$
(Planck 2018) differ by $\approx 2.0\%$, consistent with the
$1.35\%$ prime-lattice correction applied to both sectors.
\end{corollary}

\begin{remark}[OP~1/OP~2 (T4) resolution]
\label{rem:op1-tribonacci}
Theorem~\ref{thm:tribonacci} resolves OP~1 and OP~2~(T4)
\emph{in the continuum limit}:
\begin{enumerate}[label=(\roman*)]
  \item The frame factor $\gL = \mathbb{T}$ (tribonacci) is
    the unique solution of the self-referential condition
    $\gL = 1/\bsym(\gL)$ imposed by the $\mathbb{Z}/3\mathbb{Z}$
    triolic structure.
  \item The $1.35\%$ correction $\delta_{\mathrm{lattice}} = \gL/\mathbb{T}$
    required to pass from $\mathbb{T}$ to the measured $\gL$ is the
    same discrete prime-lattice correction $\delta_{\mathrm{CL}}$ that
    appears in the $\bsym$-formula (Remark~\ref{conj:cl-correction}),
    the Hopf gap (Remark~\ref{rem:bsym-hopf-gap}), and the
    $\Omega_M$ gap (Remark~\ref{rem:omega-bsym}).
  \item The full derivation of $\delta_{\mathrm{lattice}}$ from the
    prime lattice $\Gamma_P$ remains the one surviving open step
    (previously OP~1; previously OP~2~(T4)).
\end{enumerate}
\end{remark}

\begin{theorem}[First machine-fused coupled OP-entry theorem]
\label{thm:first-machine-fused-coupled-op-entry}
Assume:
\begin{enumerate}
    \item the corrected prime-lattice reading applies on a transfer-consistent closure-phase window;
    \item the first proxy-to-residual transfer theorem applies, so that the active inherited burden is read through
    \[
      \bigl(\Delta_{\mathrm{CL}},\Delta_{\Theta R}(T)\bigr)
    \]
    with $\rho_A$ as admissibility monitor;
    \item the coupled prime-lattice correction datum
    \[
      \mathfrak{C}_{\Gamma_P}^{\mathrm{cpl}}
      =
      \bigl(\Delta_{\mathrm{CL}},\Delta_{\Theta R}(T),\rho_A\bigr)
    \]
    is QAM-encodable and carried by a machine-reachable state cell;
    \item coupled prime-lattice descent is realizable in the QAM/MML/MMA layer as a finite symbolic descent/test/branch task;
    \item the closure-phase trigger slice is machine-testable by finitely checkable burden and admissibility conditions.
\end{enumerate}
Then the first post-v3.31.0 OP-entry problem admits a machine-fused coupled realization.

More precisely:
\begin{enumerate}
    \item the operative entry object is the coupled datum
    \[
      \bigl(\Delta_{\mathrm{CL}},\Delta_{\Theta R}(T),\rho_A\bigr),
    \]
    not the scalar shadow $\delta_{\mathrm{CL}}$ alone;
    \item every admissible local prime-lattice descent step can be read simultaneously
    \begin{enumerate}
        \item geometrically, as a coupled closure-phase descent step, and
        \item computationally, as a finite realizable machine step on the encoded state object;
    \end{enumerate}
    \item the closure-phase trigger slice is branch-decidable in the machine layer;
    \item therefore, whenever a finite admissible descent chain reaches the closure-phase trigger slice while preserving admissibility of $\rho_A$, immediate OP-entry follows.
\end{enumerate}
\end{theorem}

\begin{proof}
By the corrected prime-lattice reading, the scalar continuum-limit correction is only the closure-side shadow of the coupled datum
\[
  \mathfrak{C}_{\Gamma_P}^{\mathrm{cpl}}
  =
  \bigl(\Delta_{\mathrm{CL}},\Delta_{\Theta R}(T),\rho_A\bigr).
\]
By the proxy-to-residual transfer theorem, the active inherited burden is reduced to the pair
\[
  \bigl(\Delta_{\mathrm{CL}},\Delta_{\Theta R}(T)\bigr)
\]
under admissible $\rho_A$.
By QAM-encodability, this coupled datum is a machine-reachable state object.
By finite realizability of coupled prime-lattice descent, each admissible local descent is executable as a finite symbolic task in the QAM/MML/MMA layer.
By machine-testability of the closure-phase trigger slice, the machine can decide whether the updated coupled datum has reached entry, must continue descent, or has become inadmissible.
Hence the first post-v3.31.0 OP-entry problem is machine-fused: it is simultaneously a coupled geometric descent problem and a finite machine-controlled descent/test/branch problem.
Whenever the finite admissible descent chain reaches the closure-phase trigger slice while preserving $\rho_A$ admissibility, immediate OP-entry follows by the prime-lattice coupled entry lemma.
\end{proof}

\begin{corollary}[First post-v3.31.0 freeze threshold]
\label{cor:first-post-v3310-freeze-threshold}
The first machine-fused coupled OP-entry theorem marks the first natural mathematical freeze threshold for a post-v3.31.0 minor release: the line now carries a single fused claim linking the corrected prime-lattice entry geometry with the QAM/MML/MMA machine layer.
\end{corollary}

\begin{proof}
Immediate from Theorem~\ref{thm:first-machine-fused-coupled-op-entry}.
\end{proof}


\begin{definition}[Machine-fused local entry state]
\label{def:machine-fused-local-entry-state}
A \emph{machine-fused local entry state} is a tuple
\[
\mathfrak{E}_{\mathrm{loc}}
=
\bigl(
\mathbf{X}_{\Gamma_P},
\mathcal{C}_{\mathrm{mach}},
\mathfrak{D}_{\mathrm{lead}},
\mathfrak{S}_{\mathrm{trig}}
\bigr),
\]
where:
\begin{enumerate}
    \item $\mathbf{X}_{\Gamma_P}$ is the current prime-lattice machine object
    \[
    \mathbf{X}_{\Gamma_P}
    =
    \operatorname{Enc}_{\mathrm{QAM}}
    \bigl(
    \Delta_{\mathrm{CL}},
    \Delta_{\Theta R}(T),
    \rho_A
    \bigr);
    \]
    \item $\mathcal{C}_{\mathrm{mach}}$ is the current finite machine configuration carrying $\mathbf{X}_{\Gamma_P}$;
    \item $\mathfrak{D}_{\mathrm{lead}}$ is the current lead-direction tag;
    \item $\mathfrak{S}_{\mathrm{trig}}$ is the currently active trigger slice candidate.
\end{enumerate}
\end{definition}

\begin{definition}[Local machine-fused descent/test/branch operator]
\label{def:local-machine-fused-descent-test-branch-operator}
The \emph{local machine-fused descent/test/branch operator} is the composite step
\[
\mathcal{O}_{\mathrm{DTB}}
=
\mathcal{B}_{\Gamma_P}
\circ
\mathcal{T}_{\Gamma_P}
\circ
\mathcal{D}_{\Gamma_P},
\]
where:
\begin{enumerate}
    \item $\mathcal{D}_{\Gamma_P}$ is a coupled prime-lattice descent realization;
    \item $\mathcal{T}_{\Gamma_P}$ is the prime-lattice trigger test;
    \item $\mathcal{B}_{\Gamma_P}$ is the prime-lattice branch rule.
\end{enumerate}
\end{definition}

\begin{remark}[Operational reading]
\label{rem:operational-reading-local-operator}
This is the first explicit post-v3.32.0 local controller:
first lower the coupled burden pair if possible,
then test the trigger slice,
then branch into entry, continuation, or obstruction.
\end{remark}

\begin{definition}[Local obstruction certificate]
\label{def:local-obstruction-certificate}
A \emph{local obstruction certificate} is a finite machine-readable object
\[
\mathfrak{O}_{\mathrm{loc}}
\]
issued when the local machine-fused descent/test/branch operator returns
\[
\mathtt{inadmissible}
\quad\text{or}\quad
\mathtt{stalled},
\]
together with a finite explanation of which of the following failed:
\begin{enumerate}
    \item closure descent,
    \item Theta descent,
    \item admissibility of $\rho_A$,
    \item persistence of the closure-phase trigger slice,
    \item route-faithful machine realizability.
\end{enumerate}
\end{definition}

\begin{definition}[Local continuation certificate]
\label{def:local-continuation-certificate}
A \emph{local continuation certificate} is a finite machine-readable object
\[
\mathfrak{K}_{\mathrm{cont}}
\]
issued when the local operator returns
\[
\mathtt{continue},
\]
certifying that:
\begin{enumerate}
    \item the coupled descent remains admissible;
    \item the current state is still outside immediate OP-entry;
    \item at least one monotone local burden measure has strictly improved.
\end{enumerate}
\end{definition}

\begin{definition}[Local entry certificate]
\label{def:local-entry-certificate}
A \emph{local entry certificate} is a finite machine-readable object
\[
\mathfrak{K}_{\mathrm{entry}}
\]
issued when the local operator returns
\[
\mathtt{entry},
\]
certifying that the current machine-fused prime-lattice datum has entered the active closure-phase trigger slice and is immediately OP-entry certified.
\end{definition}

\begin{definition}[Machine-fused local entry protocol]
\label{def:machine-fused-local-entry-protocol}
The \emph{machine-fused local entry protocol} is the iterative rule:
\[
\mathfrak{E}_{\mathrm{loc},0}
\longmapsto
\mathfrak{E}_{\mathrm{loc},1}
\longmapsto
\mathfrak{E}_{\mathrm{loc},2}
\longmapsto \cdots
\]
obtained by repeated application of the local descent/test/branch operator, stopping when either
\[
\mathfrak{K}_{\mathrm{entry}}
\quad\text{or}\quad
\mathfrak{O}_{\mathrm{loc}}
\]
is issued.
\end{definition}

\begin{remark}[Why this is the right next step]
\label{rem:why-right-next-step-local-protocol}
The post-v3.32.0 line should no longer remain at the level of existence claims alone.
It should now either produce:
\begin{enumerate}
    \item a local entry certificate, or
    \item a local obstruction certificate.
\end{enumerate}
That is the first genuinely operational OP-entry phase.
\end{remark}

\begin{definition}[Local burden descent norm]
\label{def:local-burden-descent-norm}
A \emph{local burden descent norm} is any monotone function
\[
\mathcal{N}_{\mathrm{loc}}
\bigl(
\Delta_{\mathrm{CL}},
\Delta_{\Theta R}(T),
\rho_A
\bigr)
\]
such that:
\begin{enumerate}
    \item it decreases when both transferred burdens decrease and admissibility is preserved;
    \item it is defined only on admissible machine-fused local entry states;
    \item lower values indicate closer approach to immediate OP-entry.
\end{enumerate}
\end{definition}

\begin{proposition}[Continuation implies local progress]
\label{prop:continuation-implies-local-progress}
Assume the machine-fused local entry protocol returns a continuation certificate
\[
\mathfrak{K}_{\mathrm{cont}}
\]
at some step.
Then there exists a local burden descent norm
\[
\mathcal{N}_{\mathrm{loc}}
\]
for which the updated local entry state has strictly smaller value than the previous one.
\end{proposition}

\begin{proof}
By definition of the continuation certificate, coupled descent remains admissible, immediate OP-entry has not yet been reached, and at least one monotone local burden measure has strictly improved.
Such a measure may be taken as a local burden descent norm, yielding strict local progress.
\end{proof}

\begin{lemma}[Finite local entry-or-obstruction lemma]
\label{lem:finite-local-entry-or-obstruction-lemma}
Assume:
\begin{enumerate}
    \item the machine-fused local entry protocol is defined on the current local state;
    \item every continuation step strictly lowers a local burden descent norm;
    \item the protocol cannot descend indefinitely inside the same admissible slice without either entering the trigger slice or violating admissibility.
\end{enumerate}
Then a finite local protocol run ends with either:
\begin{enumerate}
    \item a local entry certificate, or
    \item a local obstruction certificate.
\end{enumerate}
\end{lemma}

\begin{proof}
Each continuation step strictly lowers the local burden descent norm.
By assumption, indefinite descent in the same admissible slice is impossible.
Hence the protocol must terminate after finitely many steps, either by entering the trigger slice or by losing admissibility/persistence and issuing an obstruction certificate.
\end{proof}

\begin{theorem}[First machine-fused local OP-entry protocol theorem]
\label{thm:first-machine-fused-local-op-entry-protocol-theorem}
Assume:
\begin{enumerate}
    \item the corrected prime-lattice reading applies;
    \item the first machine-fused coupled OP-entry theorem applies;
    \item the closure-phase trigger slice is machine-testable;
    \item the current line admits a machine-fused local entry state.
\end{enumerate}
Then the post-v3.32.0 line supports a finite local machine-fused OP-entry protocol.

More precisely, every admissible run of the local protocol yields after finitely many steps either:
\begin{enumerate}
    \item a local entry certificate witnessing immediate OP-entry, or
    \item a local obstruction certificate isolating the first local failure mode.
\end{enumerate}
\end{theorem}

\begin{proof}
By the first machine-fused coupled OP-entry theorem, coupled descent, trigger testing, and branch handling are already fused in the current line. Hence the local descent/test/branch operator is defined on the current machine-fused local entry state.
By the finite local entry-or-obstruction lemma, every admissible local run terminates after finitely many steps with either entry or obstruction.
\end{proof}

\begin{corollary}[Post-v3.32.0 working rule]
\label{cor:post-v3320-working-rule}
The next working line after v3.32.0 should no longer ask only whether the coupled OP-entry theorem exists abstractly.
It should explicitly run the local machine-fused entry protocol and record either:
\begin{enumerate}
    \item the first local entry certificate, or
    \item the first local obstruction certificate.
\end{enumerate}
\end{corollary}

\begin{proof}
Immediate from Theorem~\ref{thm:first-machine-fused-local-op-entry-protocol-theorem}.
\end{proof}



\begin{definition}[Default first local outcome hypothesis]
\label{def:default-first-local-outcome-hypothesis}
The \emph{default first local outcome hypothesis} for the post-v3.32.0 line is the hypothesis that the first honest terminal output of the machine-fused local entry protocol is a local obstruction certificate rather than a local entry certificate.
\end{definition}

\begin{remark}[Why obstruction is the current default]
\label{rem:why-obstruction-is-the-current-default}
At the present stage, the coupled descent mechanism is already structurally available, but the active closure-phase trigger slice is still only symbolically specified and not yet fully calibrated by explicit local threshold data. Therefore the most likely first terminal failure is not loss of the whole machine-fused architecture, but failure to certify persistence or actual reach of the trigger slice before immediate entry is declared.
\end{remark}

\begin{definition}[Trigger-persistence obstruction]
\label{def:trigger-persistence-obstruction}
A \emph{trigger-persistence obstruction} is a local obstruction certificate
\[
\mathfrak{O}_{\mathrm{loc}}^{\mathrm{trig}}
\]
issued when:
\begin{enumerate}
    \item coupled descent remains machine-realizable;
    \item the packet-side monitor remains admissible;
    \item but persistence or actual reach of the closure-phase trigger slice cannot yet be certified on the current local window.
\end{enumerate}
\end{definition}

\begin{definition}[Entry-prior versus obstruction-prior local regime]
\label{def:entry-prior-versus-obstruction-prior-local-regime}
A machine-fused local entry state is called:
\begin{enumerate}
    \item \emph{entry-prior} if the active closure-phase trigger slice is sufficiently explicit and stable that the next finite admissible local run is expected to end in a local entry certificate;
    \item \emph{obstruction-prior} if coupled descent is available but the trigger slice remains insufficiently explicit, insufficiently persistent, or insufficiently calibrated for immediate entry certification.
\end{enumerate}
\end{definition}

\begin{proposition}[The current local state is more naturally obstruction-prior than entry-prior]
\label{prop:current-local-state-more-naturally-obstruction-prior}
Assume the post-v3.32.0 line already supports:
\begin{enumerate}
    \item corrected prime-lattice reading;
    \item proxy-to-residual transfer;
    \item machine-fused descent/test/branch control.
\end{enumerate}
Assume further that the active closure-phase trigger slice is still only symbolically described and not yet equipped with fully explicit local threshold calibration.
Then the current machine-fused local entry state is more naturally obstruction-prior than entry-prior.
\end{proposition}

\begin{proof}
Under the stated assumptions, coupled descent and testing already exist as local machine operations, so the current line is not blocked at the level of basic realizability.
However, if the closure-phase trigger slice is still only symbolically described and lacks fully explicit local threshold calibration, then immediate entry cannot yet be certified purely from the present data.
Hence the more natural first terminal local outcome is an obstruction certificate isolating trigger-persistence or trigger-calibration failure rather than an entry certificate.
\end{proof}

\begin{theorem}[First local obstruction-priority theorem]
\label{thm:first-local-obstruction-priority-theorem}
Assume:
\begin{enumerate}
    \item the first machine-fused local OP-entry protocol theorem applies;
    \item the current line admits machine-fused coupled descent steps;
    \item the packet-side monitor remains admissible on the first relevant local window;
    \item the closure-phase trigger slice is not yet fully explicit or fully calibrated on that same window.
\end{enumerate}
Then the current local machine-fused state is obstruction-prior.
More precisely, the first honest terminal local outcome is expected to be a trigger-persistence obstruction
\[
\mathfrak{O}_{\mathrm{loc}}^{\mathrm{trig}},
\]
rather than an immediate local entry certificate.
\end{theorem}

\begin{proof}
By assumption, the local machine-fused protocol exists and coupled descent is available.
By admissibility of the packet-side monitor, the protocol is not forced to fail immediately by packet instability.
If, however, the closure-phase trigger slice is not yet fully explicit or calibrated, then the protocol cannot honestly certify immediate entry on the present local window.
The most natural terminal outcome is therefore a local obstruction certificate isolating the missing trigger-persistence or trigger-calibration ingredient.
\end{proof}

\begin{corollary}[First post-v3.32.0 local audit rule]
\label{cor:first-post-v3320-local-audit-rule}
The next local audit after v3.32.0 should not first ask whether immediate entry already occurs.
It should first ask which trigger-persistence or trigger-calibration ingredient is missing from the current local window, i.e. it should aim first at the first trigger-persistence obstruction certificate.
\end{corollary}

\begin{proof}
Immediate from Theorem~\ref{thm:first-local-obstruction-priority-theorem}.
\end{proof}


\begin{definition}[Trigger-calibration obstruction]
\label{def:trigger-calibration-obstruction}
A \emph{trigger-calibration obstruction} is a local obstruction certificate
\[
\mathfrak{O}_{\mathrm{loc}}^{\mathrm{cal}}
\]
issued when:
\begin{enumerate}
    \item the coupled descent/test/branch architecture is available;
    \item the active closure-phase trigger slice is recognized symbolically;
    \item but no sufficiently explicit local burden thresholds or effective local trigger inequalities are yet available to certify entry.
\end{enumerate}
\end{definition}

\begin{definition}[Trigger-persistence obstruction --- strengthened reading]
\label{def:trigger-persistence-obstruction-strengthened-reading}
A \emph{strengthened trigger-persistence obstruction} is a local obstruction certificate
\[
\mathfrak{O}_{\mathrm{loc}}^{\mathrm{pers}}
\]
issued when:
\begin{enumerate}
    \item the coupled descent/test/branch architecture is available;
    \item a symbolic closure-phase trigger slice is available;
    \item some local trigger-improving descent steps are visible;
    \item but those steps do not yet certify persistence of the trigger slice on a sufficiently long local window.
\end{enumerate}
\end{definition}

\begin{definition}[Machine-side decodability obstruction]
\label{def:machine-side-decodability-obstruction}
A \emph{machine-side decodability obstruction} is a local obstruction certificate
\[
\mathfrak{O}_{\mathrm{loc}}^{\mathrm{dec}}
\]
issued when:
\begin{enumerate}
    \item the corrected prime-lattice datum is present;
    \item but the corresponding trigger-relevant condition cannot yet be realized as a finite machine-decidable test on the active QAM/MML/MMA state.
\end{enumerate}
\end{definition}

\begin{remark}[The three natural first obstruction types]
\label{rem:the-three-natural-first-obstruction-types}
At the present stage, the first local obstruction should naturally be expected to belong to exactly one of three classes:
\begin{enumerate}
    \item trigger-calibration obstruction;
    \item trigger-persistence obstruction;
    \item machine-side decodability obstruction.
\end{enumerate}
These correspond respectively to missing local thresholds, missing stable trigger persistence, or missing test-level realizability.
\end{remark}

\begin{proposition}[Machine-side decodability is not the primary current obstruction]
\label{prop:machine-side-decodability-not-primary-current-obstruction}
Assume:
\begin{enumerate}
    \item the corrected prime-lattice reading applies;
    \item the prime-lattice machine realization theorem applies;
    \item the machine-fused local entry protocol theorem applies.
\end{enumerate}
Then a machine-side decodability obstruction is not the most natural first local obstruction of the current line.
\end{proposition}

\begin{proof}
Under the stated assumptions, the coupled prime-lattice datum is already QAM-encodable, coupled descent is MML/MMA-realizable, and the trigger problem is already machine-fused at the protocol level.
Hence the current line is not most naturally blocked by complete absence of machine-side realizability.
The more natural remaining failure modes lie at the level of trigger calibration or trigger persistence.
\end{proof}

\begin{proposition}[Current obstruction split is calibration/persistence-dominant]
\label{prop:current-obstruction-split-is-calibration-persistence-dominant}
Assume the hypotheses of Proposition~\ref{prop:machine-side-decodability-not-primary-current-obstruction}.
Assume further that the active closure-phase trigger slice is already symbolically available but not yet fully quantified by explicit local thresholds and not yet stable on a sufficiently long local window.
Then the first local obstruction is more naturally of calibration/persistence type than of decodability type.
\end{proposition}

\begin{proof}
If the trigger slice is already symbolically available, then total trigger absence is not the main issue.
If it is not yet fully quantified by explicit thresholds and not yet stable on a sufficiently long local window, then the missing ingredient lies in calibration and persistence.
By Proposition~\ref{prop:machine-side-decodability-not-primary-current-obstruction}, this is more natural than a primary machine-side decodability failure.
\end{proof}

\begin{definition}[Current dominant first local obstruction class]
\label{def:current-dominant-first-local-obstruction-class}
The \emph{current dominant first local obstruction class} is the class
\[
\mathfrak{O}_{\mathrm{loc}}^{\ast}
\in
\{
\mathfrak{O}_{\mathrm{loc}}^{\mathrm{cal}},
\mathfrak{O}_{\mathrm{loc}}^{\mathrm{pers}},
\mathfrak{O}_{\mathrm{loc}}^{\mathrm{dec}}
\}
\]
which is most naturally expected to occur first on the current local window.
\end{definition}

\begin{theorem}[First local obstruction classification theorem]
\label{thm:first-local-obstruction-classification-theorem}
Assume:
\begin{enumerate}
    \item the first local obstruction-priority theorem applies;
    \item the corrected prime-lattice reading applies;
    \item the prime-lattice machine realization theorem applies;
    \item the closure-phase trigger slice is already symbolically available;
    \item the trigger slice is not yet fully locally calibrated and not yet certified persistent on a sufficiently long local window.
\end{enumerate}
Then the current dominant first local obstruction class is not machine-side decodability but calibration/persistence type.

More precisely,
\[
\mathfrak{O}_{\mathrm{loc}}^{\ast}
\in
\{
\mathfrak{O}_{\mathrm{loc}}^{\mathrm{cal}},
\mathfrak{O}_{\mathrm{loc}}^{\mathrm{pers}}
\},
\]
and the current line should be read as calibration/persistence-obstruction-prior.
\end{theorem}

\begin{proof}
By the prime-lattice machine realization theorem, the coupled datum is already machine-encodable and machine-fused descent/test/branch control is available, so pure decodability failure is not the most natural first obstruction.
Because the closure-phase trigger slice is already symbolically available but not yet fully calibrated or persistent on a sufficiently long local window, the remaining failure modes lie in trigger calibration and trigger persistence.
Hence the current dominant obstruction class is of calibration/persistence type.
\end{proof}

\begin{corollary}[Next post-v3.32.0 local audit question]
\label{cor:next-post-v3320-local-audit-question}
The next local audit after the current r13 stage should first ask:
\begin{enumerate}
    \item what explicit local trigger inequalities are still missing;
    \item and what minimum persistence window is still missing
\end{enumerate}
before treating machine-side decodability as the primary obstruction.
\end{corollary}

\begin{proof}
Immediate from Theorem~\ref{thm:first-local-obstruction-classification-theorem}.
\end{proof}

\begin{definition}[Local closure-phase trigger inequality schema]
\label{def:local-closure-phase-trigger-inequality-schema}
A \emph{local closure-phase trigger inequality schema} is a finite family of inequalities of the form
\[
\Delta_{\mathrm{CL}} \le \varepsilon_{\mathrm{CL}},
\qquad
\Delta_{\Theta R}(T) \le \varepsilon_{\Theta},
\qquad
\rho_A \in \mathfrak{A}_{\mathrm{adm}}^{\mathrm{safe}},
\]
together with the side condition that the current transition remains closure-phase led on the same local window.
Here
\[
\varepsilon_{\mathrm{CL}}>0,
\qquad
\varepsilon_{\Theta}>0
\]
are local trigger tolerances and
\[
\mathfrak{A}_{\mathrm{adm}}^{\mathrm{safe}}\subseteq \mathfrak{A}_{\mathrm{adm}}
\]
is a distinguished safe admissibility region for the packet-side monitor.
\end{definition}

\begin{remark}[Meaning of the inequality schema]
\label{rem:meaning-of-the-inequality-schema}
The current line does not yet require a final closed-form trigger formula.
At the present stage, what is missing is already captured by the weaker data of explicit local tolerances for the coupled burdens together with a safe admissibility band for \(\rho_A\).
\end{remark}

\begin{definition}[Local closure-phase persistence window of length \texorpdfstring{$m$}{m}]
\label{def:local-closure-phase-persistence-window-of-length-m}
Let \(m\ge 1\).
A \emph{local closure-phase persistence window of length \(m\)} is a block of successive local machine-fused protocol steps
\[
\mathfrak{E}_{\mathrm{loc},n},\mathfrak{E}_{\mathrm{loc},n+1},\dots,\mathfrak{E}_{\mathrm{loc},n+m-1}
\]
such that throughout the block:
\begin{enumerate}
    \item the lead direction remains closure-phase led;
    \item the trigger inequality schema is satisfied monotonically or remains stable once entered;
    \item the packet-side monitor stays in \(\mathfrak{A}_{\mathrm{adm}}^{\mathrm{safe}}\);
    \item the same route-faithful local slice remains active.
\end{enumerate}
\end{definition}

\begin{definition}[Minimum local trigger-persistence length]
\label{def:minimum-local-trigger-persistence-length}
A positive integer
\[
m_{\min}\ge 1
\]
is called a \emph{minimum local trigger-persistence length} if immediate OP-entry is not certified merely by entering the local trigger inequality schema once, but only after the schema persists on a closure-phase persistence window of length at least \(m_{\min}\).
\end{definition}

\begin{remark}[Why a persistence length is needed]
\label{rem:why-a-persistence-length-is-needed}
The current obstruction analysis suggests that a one-step near-entry fluctuation is not yet trustworthy.
What is missing is not just smallness of the coupled burdens, but stable smallness on a sufficiently long local window.
\end{remark}

\begin{definition}[Trigger-calibration data package]
\label{def:trigger-calibration-data-package}
The \emph{trigger-calibration data package} is the tuple
\[
\mathfrak{P}_{\mathrm{trig}}^{\mathrm{loc}}
:=
\bigl(
\varepsilon_{\mathrm{CL}},
\varepsilon_{\Theta},
\mathfrak{A}_{\mathrm{adm}}^{\mathrm{safe}},
 m_{\min}
\bigr)
\]
consisting of local burden tolerances, a safe admissibility region, and a minimum persistence length.
\end{definition}

\begin{proposition}[What the current obstruction is really asking for]
\label{prop:what-the-current-obstruction-is-really-asking-for}
Assume the first local obstruction classification theorem.
Then the current calibration/persistence obstruction is equivalent to the absence of a sufficiently explicit trigger-calibration data package.
\end{proposition}

\begin{proof}
By the obstruction classification theorem, the current obstruction is not primarily one of machine-side decodability, but of calibration/persistence type.
Calibration requires explicit local burden tolerances and a safe admissibility region, while persistence requires a minimum stable window length.
Therefore the current obstruction is exactly the absence of a sufficiently explicit trigger-calibration data package.
\end{proof}

\begin{definition}[Locally calibrated closure-phase trigger slice]
\label{def:locally-calibrated-closure-phase-trigger-slice}
A closure-phase trigger slice is called \emph{locally calibrated} if it is equipped with a trigger-calibration data package
\[
\mathfrak{P}_{\mathrm{trig}}^{\mathrm{loc}}
\]
and immediate OP-entry is certified whenever the associated local closure-phase trigger inequality schema persists for at least \(m_{\min}\) consecutive local protocol steps.
\end{definition}

\begin{lemma}[First local calibration-to-entry lemma]
\label{lem:first-local-calibration-to-entry-lemma}
Assume:
\begin{enumerate}
    \item the machine-fused local entry protocol theorem applies;
    \item a locally calibrated closure-phase trigger slice is available;
    \item the current local protocol run enters the associated trigger inequality schema and remains there for at least \(m_{\min}\) consecutive steps.
\end{enumerate}
Then the local protocol yields a local entry certificate.
\end{lemma}

\begin{proof}
By local calibration, entry is certified exactly when the trigger inequality schema persists on a closure-phase window of length at least \(m_{\min}\).
If the current local protocol run satisfies this condition, then by definition the corresponding state is immediately OP-entry certified and the protocol returns a local entry certificate.
\end{proof}

\begin{theorem}[First local trigger calibration theorem]
\label{thm:first-local-trigger-calibration-theorem}
Assume:
\begin{enumerate}
    \item the first local obstruction classification theorem applies;
    \item the machine-fused local entry protocol theorem applies;
    \item the active obstruction is of calibration/persistence type.
\end{enumerate}
Then the next missing local ingredient is exactly the explicit trigger-calibration data package
\[
\mathfrak{P}_{\mathrm{trig}}^{\mathrm{loc}}
=
\bigl(
\varepsilon_{\mathrm{CL}},
\varepsilon_{\Theta},
\mathfrak{A}_{\mathrm{adm}}^{\mathrm{safe}},
 m_{\min}
\bigr).
\]
Moreover, once such a package is fixed, the remaining local question reduces to whether the current protocol run can enter and maintain the associated inequality schema for at least \(m_{\min}\) consecutive steps.
\end{theorem}

\begin{proof}
By the first local obstruction classification theorem, the dominant current obstruction is not decodability but calibration/persistence type.
By Proposition~\ref{prop:what-the-current-obstruction-is-really-asking-for}, this is exactly the absence of explicit local burden tolerances, a safe admissibility region, and a minimum persistence length.
These data are precisely the trigger-calibration data package.
Once the package is fixed, the only remaining local issue is whether the current machine-fused protocol run satisfies the corresponding persistence requirement.
\end{proof}

\begin{corollary}[Next post-r14 working rule]
\label{cor:next-post-r14-working-rule}
The next working step after the current obstruction analysis should not try to prove immediate OP-entry in one jump.
It should first determine:
\begin{enumerate}
    \item plausible local values or bounds for \(\varepsilon_{\mathrm{CL}}\) and \(\varepsilon_{\Theta}\);
    \item a safe admissibility subregion \(\mathfrak{A}_{\mathrm{adm}}^{\mathrm{safe}}\);
    \item a first plausible minimum persistence length \(m_{\min}\).
\end{enumerate}
\end{corollary}

\begin{proof}
Immediate from Theorem~\ref{thm:first-local-trigger-calibration-theorem}.
\end{proof}



\begin{definition}[Default local trigger orientation]
\label{def:default-local-trigger-orientation}
The \emph{default local trigger orientation} is the working hypothesis that the first explicit local closure-phase trigger calibration should satisfy:
\begin{enumerate}
    \item the closure tolerance is the stricter burden tolerance,
    \[
    \varepsilon_{\mathrm{CL}} \le \varepsilon_{\Theta};
    \]
    \item the safe admissibility region for \(\rho_A\) is taken as a narrow stability band around the current admissible packet-side monitor regime;
    \item the minimum persistence length is not read as one-step only, but as a genuinely multi-step window
    \[
    m_{\min}>1.
    \]
\end{enumerate}
\end{definition}

\begin{remark}[Why closure tolerance should be stricter first]
\label{rem:why-closure-tolerance-should-be-stricter-first}
The current obstruction analysis suggests that the coupled local burden is not blocked first by lack of Theta-side machine decodability, but by insufficiently stabilized closure-phase persistence.
Hence the first local calibration should be closure-tight rather than Theta-tight.
\end{remark}

\begin{definition}[First plausible default trigger package]
\label{def:first-plausible-default-trigger-package}
The \emph{first plausible default trigger package} is any trigger-calibration data package
\[
\mathfrak{P}_{\mathrm{trig}}^{\mathrm{loc,def}}
:=
\bigl(
\varepsilon_{\mathrm{CL}}^{\mathrm{def}},
\varepsilon_{\Theta}^{\mathrm{def}},
\mathfrak{A}_{\mathrm{adm}}^{\mathrm{safe,def}},
 m_{\min}^{\mathrm{def}}
\bigr)
\]
satisfying:
\begin{enumerate}
    \item
    \[
    0<\varepsilon_{\mathrm{CL}}^{\mathrm{def}} \le \varepsilon_{\Theta}^{\mathrm{def}};
    \]
    \item
    \[
    \mathfrak{A}_{\mathrm{adm}}^{\mathrm{safe,def}}
    \subseteq
    \mathfrak{A}_{\mathrm{adm}};
    \]
    \item
    \[
    m_{\min}^{\mathrm{def}}\ge 2.
    \]
\end{enumerate}
\end{definition}

\begin{remark}[Interpretation of the default package]
\label{rem:interpretation-of-the-default-package}
This package does not yet assign final numerical values.
It only fixes the first local reading:
closure burden must enter a stricter small-burden regime than the routed Theta burden, and one-step near-entry is not yet trusted as a sufficient certificate.
\end{remark}

\begin{proposition}[Default local trigger package is obstruction-compatible]
\label{prop:default-local-trigger-package-is-obstruction-compatible}
Assume the first local obstruction classification theorem.
Then the first plausible default trigger package is compatible with the current obstruction reading.
\end{proposition}

\begin{proof}
The current obstruction has been classified as calibration/persistence type rather than machine-side decodability type.
A stricter closure tolerance together with a safe admissibility subregion and a persistence length at least two directly addresses that obstruction pattern: it stabilizes the closure burden first, keeps the packet-side monitor conservative, and excludes one-step transient near-entry events.
Hence the default package is compatible with the current obstruction reading.
\end{proof}

\begin{proposition}[Why one-step persistence should not be the default]
\label{prop:why-one-step-persistence-should-not-be-the-default}
Under the present obstruction-prior reading, the choice
\[
m_{\min}=1
\]
should not be the first default calibration.
\end{proposition}

\begin{proof}
The obstruction analysis explicitly identified trigger-persistence instability as part of the first local failure mode.
If one were to set \(m_{\min}=1\), then a single transient entrance into the local trigger inequality schema would already be counted as sufficient, which removes exactly the persistence safeguard that the current obstruction analysis says is missing.
Therefore one-step persistence is not the correct first default choice.
\end{proof}

\begin{theorem}[First default local trigger calibration theorem]
\label{thm:first-default-local-trigger-calibration-theorem}
Assume:
\begin{enumerate}
    \item the first local trigger calibration theorem applies;
    \item the current obstruction is of calibration/persistence type;
    \item the closure-phase-first local reading remains active.
\end{enumerate}
Then the first reasonable local default calibration should be chosen from a first plausible default trigger package
\[
\mathfrak{P}_{\mathrm{trig}}^{\mathrm{loc,def}}
=
\bigl(
\varepsilon_{\mathrm{CL}}^{\mathrm{def}},
\varepsilon_{\Theta}^{\mathrm{def}},
\mathfrak{A}_{\mathrm{adm}}^{\mathrm{safe,def}},
 m_{\min}^{\mathrm{def}}
\bigr)
\]
with
\[
0<\varepsilon_{\mathrm{CL}}^{\mathrm{def}} \le \varepsilon_{\Theta}^{\mathrm{def}},
\qquad
m_{\min}^{\mathrm{def}}\ge 2.
\]
In particular, the first local trigger calibration is closure-prior and persistence-prior rather than one-step permissive.
\end{theorem}

\begin{proof}
By the first local trigger calibration theorem, what is missing is an explicit trigger-calibration package.
By the obstruction classification, the leading obstruction is calibration/persistence type and not machine-side decodability.
Therefore the first calibration should tighten the closure burden first, allow at least as much tolerance on the routed Theta burden, keep the packet-side admissibility band conservative, and require a genuinely persistent window rather than a one-step fluctuation.
These conditions are exactly those of the first plausible default trigger package.
\end{proof}

\begin{corollary}[Current post-r15 default working rule]
\label{cor:current-post-r15-default-working-rule}
At the present stage, the first local trigger audit should proceed as if:
\begin{enumerate}
    \item closure smallness must be at least as strict as Theta smallness;
    \item the packet-side admissibility band must be chosen conservatively;
    \item the first honest local persistence length should be read as
    \[
    m_{\min}\ge 2.
    \]
\end{enumerate}
Only after this default calibration fails should one relax toward a more permissive local trigger reading.
\end{corollary}

\begin{proof}
Immediate from Theorem~\ref{thm:first-default-local-trigger-calibration-theorem}.
\end{proof}


\begin{definition}[Discretized Teichm\"uller persistence stream]
\label{def:discretized-teichmueller-persistence-stream}
Let
\[
X_n
\]
denote the current local machine-fused state and let
\[
T_N(X_n)
\]
be the local Monte-Carlo Teichm\"uller distortion surrogate.

The \emph{discretized Teichm\"uller persistence stream} is the symbolic stream
\[
\Sigma_{\mathrm{pers}}
=
(\sigma_n)_{n\ge 0},
\]
where each symbol \(\sigma_n\) records the discretized local status of
\[
\bigl(
\Delta_{\mathrm{CL}}(n),
\Delta_{\Theta R}(T)(n),
\rho_A(n),
T_N(X_n),
T_N(X_{n+1}),
|T_N(X_{n+1})-T_N(X_n)|
\bigr)
\]
relative to the active local trigger thresholds.
\end{definition}

\begin{definition}[Kasiski persistence scan]
\label{def:kasiski-persistence-scan}
The \emph{Kasiski persistence scan} is the search for repeated symbolic blocks inside
\[
\Sigma_{\mathrm{pers}}
\]
in order to detect recurring local trigger-near, persistence-fail, or admissibility-fail patterns.
\end{definition}

\begin{remark}[Role of Kasiski]
\label{rem:role-of-kasiski}
Kasiski is used here as a repetition and spacing detector on the discretized persistence stream.
It does not decide admissibility; it only detects recurring local obstruction patterns.
\end{remark}

\begin{definition}[Hashed persistence window]
\label{def:hashed-persistence-window}
A \emph{hashed persistence window} of length \(m\) is a hash value assigned to a local window
\[
(\sigma_n,\sigma_{n+1},\dots,\sigma_{n+m-1})
\]
of the discretized Teichm\"uller persistence stream.
\end{definition}

\begin{definition}[Markov persistence surface]
\label{def:markov-persistence-surface}
The \emph{Markov persistence surface} is the transition surface obtained from the empirical transition frequencies between hashed persistence windows or their reduced regime classes.
\end{definition}

\begin{remark}[Role of Markov control]
\label{rem:role-of-markov-control}
The Markov layer is not a primary admissibility criterion.
It is only the readable transition shadow of the shell-/trigger-restricted local persistence dynamics.
\end{remark}

\begin{definition}[Newton persistence refinement]
\label{def:newton-persistence-refinement}
The \emph{Newton persistence refinement} is a local Newton-type correction step applied only after a persistence regime has already been selected and a smooth surrogate defect function has been chosen.
\end{definition}

\begin{remark}[Why Newton comes last]
\label{rem:why-newton-comes-last}
Newton is not used to discover the regime.
It is used only to sharpen the local estimate once the regime has already been isolated by Kasiski/hash/Markov filtering.
\end{remark}

\begin{definition}[Bombe-sieve persistence elimination layer]
\label{def:bombe-sieve-persistence-elimination}
The \emph{bombe-sieve persistence elimination layer} is the contradiction-sensitive elimination of local continuation menus by:
\begin{enumerate}
    \item a trigger-sensitive sieve,
    \item bombe-style rejection of locally incompatible continuation branches,
    \item optional nested backtracking across deeper local menus.
\end{enumerate}
\end{definition}

\begin{definition}[Spectral menu readout]
\label{def:spectral-menu-readout}
A \emph{spectral menu readout} is the local ranking of a shell-admissible continuation menu by coding it into a self-adjoint menu operator and reading the preferred continuation as the top eigendirection.
\end{definition}

\begin{definition}[Vernam-like persistence stream heuristic]
\label{def:vernam-like-persistence-stream-heuristic}
The \emph{Vernam-like persistence stream heuristic} is the working interpretation that the discretized Teichm\"uller persistence stream behaves like an observable stream generated by a hidden local structural source, so that repeated local obstruction patterns may be used to infer the generator rather than merely the visible symbols.
\end{definition}

\begin{remark}[How to read the Vernam-like heuristic]
\label{rem:how-to-read-the-vernam-like-heuristic}
This is not an identity claim with literal cryptographic stream ciphers.
It is a search heuristic: the visible persistence symbols are treated as a surface stream whose hidden local generator is to be reconstructed from repetition, spacing, transition structure, and contradiction-sensitive elimination.
\end{remark}

\begin{definition}[Crypto-guided persistence analysis pipeline]
\label{def:crypto-guided-persistence-analysis-pipeline}
The \emph{crypto-guided persistence analysis pipeline} is the local tool chain
\[
\text{Kasiski}
\to
\text{hash}
\to
\text{Markov}
\to
\text{Bombe/Sieve}
\to
\text{spectral menu readout}
\to
\text{Newton refinement},
\]
applied to the discretized Teichm\"uller persistence stream.
\end{definition}

\begin{theorem}[First crypto-guided persistence scan theorem]
\label{thm:first-crypto-guided-persistence-scan-theorem}
Assume:
\begin{enumerate}
    \item a local Monte-Carlo Teichm\"uller persistence stream has been constructed;
    \item the active closure/Theta/admissibility thresholds are fixed locally;
    \item admissibility remains decided by closure-compatible trigger logic.
\end{enumerate}
Then the crypto-guided persistence analysis pipeline is admissible as a subordinate persistence-analysis layer.
It may be used to detect recurring persistence-failure patterns, compress local continuation menus, and sharpen local trigger calibration without replacing the closure-compatible admissibility criterion.
\end{theorem}

\begin{proof}
The earlier tool-guided compression logic already admits Kasiski, hash/rainbow compression, Markov guidance, contradiction-sensitive elimination, spectral ranking, and Newton refinement provided they remain subordinate to structural admissibility.
Here the same discipline is applied to the discretized Teichm\"uller persistence stream.
Hence the pipeline is admissible as a subordinate persistence-analysis layer and may be used to detect, compress, rank, and refine local persistence behavior without replacing closure-compatible trigger logic itself.
\end{proof}

\begin{corollary}[Current post-r16 persistence working rule]
\label{cor:current-post-r16-persistence-working-rule}
The next local persistence audit should not rely on pointwise burden smallness alone.
It should:
\begin{enumerate}
    \item build the discretized Teichm\"uller persistence stream;
    \item scan it by Kasiski for recurring local obstruction patterns;
    \item hash and compress local windows;
    \item read the resulting regime dynamics through a Markov persistence surface;
    \item eliminate locally contradictory branches by bombe/sieve filtering;
    \item rank the surviving continuation menu spectrally;
    \item apply Newton refinement only at the final local sharpening stage.
\end{enumerate}
\end{corollary}

\begin{proof}
Immediate from Theorem~\ref{thm:first-crypto-guided-persistence-scan-theorem}.
\end{proof}



\begin{definition}[Local persistence symbols]
\label{def:local-persistence-symbols}
Let \(\Sigma_{\mathrm{pers}}=(\sigma_n)_{n\ge 0}\) be the discretized Teichm\"uller persistence stream.
The following local symbols are distinguished:
\begin{enumerate}
    \item \(\mathtt{NT}\) (near-trigger): the local burden pair satisfies the pointwise closure/Theta smallness bounds but the minimal persistence window is not yet certified;
    \item \(\mathtt{PF}\) (persistence-fail): the next local state leaves the candidate closure-phase trigger slice or fails the Teichm\"uller persistence control while remaining machine-decodable;
    \item \(\mathtt{CF}\) (calibration-fail): one of the active closure/Theta trigger inequalities is not yet met;
    \item \(\mathtt{AF}\) (admissibility-fail): the packet-side monitor leaves the safe admissibility zone;
    \item \(\mathtt{EN}\) (entry): the local state lies in the calibrated closure-phase trigger slice.
\end{enumerate}
\end{definition}

\begin{definition}[First recurring local failure motif]
\label{def:first-recurring-local-failure-motif}
The \emph{first recurring local failure motif} is the shortest repeated symbolic block
\[
\omega_{\mathrm{fail}}
=
(\sigma_n,\dots,\sigma_{n+\ell-1})
\]
with \(\ell\ge 2\) detected by the Kasiski persistence scan and interpreted as the first persistent local obstruction pattern.
\end{definition}

\begin{definition}[Near-trigger / persistence-fail motif]
\label{def:near-trigger-persistence-fail-motif}
The \emph{near-trigger / persistence-fail motif} is the symbolic block
\[
\omega_{\mathrm{NT\to PF}}
:=
(\mathtt{NT},\mathtt{PF}).
\]
It records a local run in which the first step reaches pointwise trigger-near smallness, but the second step fails to preserve the closure-phase trigger slice or the corresponding Teichm\"uller persistence control.
\end{definition}

\begin{remark}[Why this is the natural first failure candidate]
\label{rem:why-natural-first-failure-candidate}
The current post-v3.32.0 line is obstruction-prior rather than entry-prior, and its first classified obstruction is of trigger calibration/persistence type rather than machine-side decodability type.
Hence the first repeated failure pattern should be sought among motifs that are already near local trigger conditions but fail on the required two-step persistence window.
\end{remark}

\begin{proposition}[Machine-side decodability is not the first recurring failure source]
\label{prop:machine-side-decodability-not-first-recurring-failure-source}
Assume the prime-lattice machine realization theorem and the first machine-fused local OP-entry protocol theorem.
Then the first recurring local failure motif is not forced by machine-side decodability failure alone.
\end{proposition}

\begin{proof}
Under the machine-fused realization hypotheses, descent, trigger testing, and branching are already available as finite machine-controlled operations on the prime-lattice machine object.
Therefore a repeated first failure pattern is not naturally attributed to absence of local machine decodability itself, but to failure of trigger calibration or persistence within the already machine-realizable regime.
\end{proof}

\begin{theorem}[First recurring local persistence-failure reading]
\label{thm:first-recurring-local-persistence-failure-reading}
Assume:
\begin{enumerate}
    \item the closure-phase-first local audit remains the preferred default reading;
    \item the default local trigger orientation with \(\varepsilon_{\mathrm{CL}}\le \varepsilon_{\Theta}\) and \(m_{\min}\ge 2\) is adopted;
    \item the current line is obstruction-prior rather than entry-prior;
    \item the crypto-guided persistence analysis pipeline is applied to the local discretized Teichm\"uller persistence stream.
\end{enumerate}
Then the first recurring local failure motif is read by default as
\[
\omega_{\mathrm{fail}}
=
\omega_{\mathrm{NT\to PF}}
=
(\mathtt{NT},\mathtt{PF}),
\]
that is: a trigger-near first step followed by failure of two-step closure-phase persistence.
\end{theorem}

\begin{proof}
By obstruction-priority, immediate entry is not the default first local outcome.
By the earlier obstruction classification, the first local obstruction is of trigger calibration/persistence type rather than machine-side decodability type.
Under the default trigger orientation and the requirement \(m_{\min}\ge 2\), the most economical local failure pattern is therefore not a one-step calibration collapse before near-trigger status is reached, but a run in which the first step attains trigger-near pointwise smallness and the second step fails to preserve the closure-phase trigger slice.
The crypto-guided persistence scan is designed exactly to detect repeated symbolic blocks of this form.
Hence the default first recurring local failure motif is \((\mathtt{NT},\mathtt{PF})\).
\end{proof}

\begin{corollary}[Current first persistence-stream reading]
\label{cor:current-first-persistence-stream-reading}
The current post-v3.32.0 persistence line should be read first through the motif
\[
(\mathtt{NT},\mathtt{PF}),
\]
rather than through immediate entry or pure decodability failure.
Equivalently, the present first local task is to identify which closure-phase persistence ingredient breaks between the near-trigger first step and the failed second step.
\end{corollary}

\begin{proof}
Immediate from Theorem~\ref{thm:first-recurring-local-persistence-failure-reading} and Proposition~\ref{prop:machine-side-decodability-not-first-recurring-failure-source}.
\end{proof}


\begin{definition}[Theta-drift persistence-fail subtype]
\label{def:theta-drift-persistence-fail-subtype}
The symbol
\[
\mathtt{PF}_{\Theta}
\]
denotes the persistence-fail subtype in which:
\begin{enumerate}
    \item the first local step is near-trigger;
    \item the second local step fails because the phase-side component loses the required closure-phase persistence control;
    \item the packet-side admissibility monitor remains inside the safe zone.
\end{enumerate}
\end{definition}

\begin{definition}[Teichm\"uller-distortion persistence-fail subtype]
\label{def:teichmueller-distortion-persistence-fail-subtype}
The symbol
\[
\mathtt{PF}_{T}
\]
denotes the persistence-fail subtype in which:
\begin{enumerate}
    \item the first local step is near-trigger;
    \item the pointwise burden inequalities remain approximately trigger-near;
    \item the second local step fails because the local Teichm\"uller distortion surrogate or its one-step variation leaves the allowed persistence band.
\end{enumerate}
\end{definition}

\begin{definition}[Admissibility-edge persistence-fail subtype]
\label{def:admissibility-edge-persistence-fail-subtype}
The symbol
\[
\mathtt{PF}_{A}
\]
denotes the persistence-fail subtype in which:
\begin{enumerate}
    \item the first local step is near-trigger;
    \item the second local step fails because the packet-side monitor \(\rho_A\) reaches or leaves the safe admissibility edge;
    \item the corresponding closure-phase persistence window therefore collapses for admissibility reasons.
\end{enumerate}
\end{definition}

\begin{remark}[How the three subtypes differ]
\label{rem:how-the-three-subtypes-differ}
The subtype \(\mathtt{PF}_{\Theta}\) is phase-led and burden-side.
The subtype \(\mathtt{PF}_{T}\) is geometric and distortion-led.
The subtype \(\mathtt{PF}_{A}\) is guard-led and packet-admissibility driven.
\end{remark}

\begin{proposition}[Admissibility-edge contact is not the default first persistence-fail subtype]
\label{prop:admissibility-edge-contact-not-default-first-pf-subtype}
Assume the current line is read through the default local trigger orientation with conservative packet-side admissibility control.
Then \(\mathtt{PF}_{A}\) is not the default first persistence-fail subtype.
\end{proposition}

\begin{proof}
The current line treats \(\rho_A\) primarily as an admissibility monitor rather than as a coequal core burden, and the active default trigger package keeps the safe admissibility band deliberately conservative.
Hence the first local failure is not naturally read as immediate edge contact of \(\rho_A\), but as failure of the closure-phase persistence mechanism before the admissibility edge is reached.
\end{proof}

\begin{proposition}[The first persistence-fail subtype is phase/geometry-prior]
\label{prop:first-persistence-fail-subtype-is-phase-geometry-prior}
Assume the closure-phase-first local audit, the default trigger orientation, and the first recurring persistence-failure reading
\[
(\mathtt{NT},\mathtt{PF}).
\]
Then the first persistence-fail subtype is read by default as phase/geometry-prior, i.e. as lying in the reduced pair
\[
\{\mathtt{PF}_{\Theta},\mathtt{PF}_{T}\},
\]
rather than as an admissibility-edge subtype \(\mathtt{PF}_{A}\).
\end{proposition}

\begin{proof}
Under the closure-phase-first reading, the active mixed persistence burden is carried first by the closure/Theta pair, while \(\rho_A\) remains a guard variable.
By Proposition~\ref{prop:admissibility-edge-contact-not-default-first-pf-subtype}, immediate admissibility-edge contact is not the default first subtype.
Hence the first persistence-fail subtype is phase/geometry-prior.
\end{proof}

\begin{theorem}[First local persistence-fail subtype theorem]
\label{thm:first-local-persistence-fail-subtype-theorem}
Assume:
\begin{enumerate}
    \item the first recurring local failure motif is
    \[
    (\mathtt{NT},\mathtt{PF});
    \]
    \item the closure-phase-first local audit remains the default reading;
    \item the packet-side admissibility monitor remains conservative on the first near-trigger step.
\end{enumerate}
Then the first local persistence-fail subtype is not read by default as admissibility-edge contact.
Instead, it is read by default as a phase/geometry-prior break:
either
\[
\mathtt{PF}_{\Theta}
\]
or
\[
\mathtt{PF}_{T},
\]
with the current line favoring the reading that Theta-side drift is the primary burden-side break and Teichm\"uller distortion is its geometric witness.
\end{theorem}

\begin{proof}
By Proposition~\ref{prop:first-persistence-fail-subtype-is-phase-geometry-prior}, the first subtype lies in the pair \(\{\mathtt{PF}_{\Theta},\mathtt{PF}_{T}\}\) rather than in \(\mathtt{PF}_{A}\).
Under the closure-phase-first reading, the active burden is phase-led before it is envelope-led, while the Teichm\"uller control records the geometric persistence of that phase-led regime.
Therefore the default local reading is that Theta-drift is the primary burden-side break and Teichm\"uller distortion is the corresponding geometric witness of the failed second step.
\end{proof}

\begin{corollary}[Current post-r18 local diagnosis rule]
\label{cor:current-post-r18-local-diagnosis-rule}
The next local diagnosis should not first ask whether \(\rho_A\) touched the admissibility edge.
It should first ask:
\begin{enumerate}
    \item did the phase-side burden drift out of the closure-phase persistence window;
    \item and if so, what Teichm\"uller distortion witness recorded that loss of persistence?
\end{enumerate}
\end{corollary}

\begin{proof}
Immediate from Theorem~\ref{thm:first-local-persistence-fail-subtype-theorem}.
\end{proof}


\begin{definition}[Local Theta-drift witness]
\label{def:local-theta-drift-witness}
A \emph{local Theta-drift witness} is a finite witness object
\[
\mathfrak{W}_{\Theta}^{\mathrm{loc}}(n)
\]
attached to a near-trigger first step and the following failed second step, certifying that:
\begin{enumerate}
    \item the first step is of type \(\mathtt{NT}\);
    \item the second step is of type \(\mathtt{PF}_{\Theta}\) or phase/geometry-prior persistence failure;
    \item the phase-side burden leaves the active closure-phase persistence window on the second step while the packet-side admissibility monitor remains inside the conservative safe zone.
\end{enumerate}
\end{definition}

\begin{definition}[Local Teichm\"uller drift witness]
\label{def:local-teichmueller-drift-witness}
A \emph{local Teichm\"uller drift witness} is a finite witness object
\[
\mathfrak{W}_{T}^{\mathrm{loc}}(n)
\]
certifying that the local Teichm\"uller distortion surrogate or its one-step variation crosses the active persistence band between the near-trigger first step and the failed second step.
\end{definition}

\begin{remark}[Burden-side versus geometric witness]
\label{rem:burden-side-versus-geometric-witness}
The Theta-drift witness and the Teichm\"uller drift witness play different roles.
The first records the burden-side loss of closure-phase persistence.
The second records the geometric distortion signature of that same loss.
\end{remark}

\begin{definition}[Local coupled drift witness package]
\label{def:local-coupled-drift-witness-package}
The \emph{local coupled drift witness package} is the pair
\[
\mathfrak{W}_{\mathrm{drift}}^{\mathrm{loc}}(n)
:=
\bigl(
\mathfrak{W}_{\Theta}^{\mathrm{loc}}(n),
\mathfrak{W}_{T}^{\mathrm{loc}}(n)
\bigr),
\]
read on the current persistence stream whenever the local failure motif
\[
(\mathtt{NT},\mathtt{PF})
\]
is refined to a phase/geometry-prior break.
\end{definition}

\begin{definition}[Local Theta-drift witness test]
\label{def:local-theta-drift-witness-test}
The \emph{local Theta-drift witness test} asks whether, for some local index \(n\):
\begin{enumerate}
    \item the first step satisfies the active near-trigger inequalities;
    \item the second step violates the phase-side persistence requirement before admissibility edge contact occurs;
    \item a Teichm\"uller distortion witness is simultaneously present or becomes newly active on that second step.
\end{enumerate}
\end{definition}

\begin{proposition}[Theta-drift witness packages refine the default persistence-fail reading]
\label{prop:theta-drift-witness-packages-refine-default-persistence-fail-reading}
Assume the first recurring local failure motif is
\[
(\mathtt{NT},\mathtt{PF})
\]
and the first persistence-fail subtype is phase/geometry-prior.
Then every successful local Theta-drift witness test yields a local coupled drift witness package
\[
\mathfrak{W}_{\mathrm{drift}}^{\mathrm{loc}}(n)
\]
refining the persistence-fail reading.
\end{proposition}

\begin{proof}
If the local Theta-drift witness test succeeds, then the failed second step is not merely an undifferentiated persistence failure: it is a phase-side drift event accompanied by a geometric distortion witness while admissibility remains conservative.
That is exactly the data recorded by the two witness objects and therefore by their paired witness package.
\end{proof}

\begin{theorem}[First local Theta-drift witness theorem]
\label{thm:first-local-theta-drift-witness-theorem}
Assume:
\begin{enumerate}
    \item the current post-v3.32.0 persistence line is read through the first recurring motif
    \[
    (\mathtt{NT},\mathtt{PF});
    \]
    \item the first persistence-fail subtype is phase/geometry-prior;
    \item the local Theta-drift witness test succeeds at some first relevant local index \(n\).
\end{enumerate}
Then the current line admits a first local Theta-drift witness package
\[
\mathfrak{W}_{\mathrm{drift}}^{\mathrm{loc}}(n)
=
\bigl(
\mathfrak{W}_{\Theta}^{\mathrm{loc}}(n),
\mathfrak{W}_{T}^{\mathrm{loc}}(n)
\bigr),
\]
with the following reading:
\begin{enumerate}
    \item the first failed second step is burden-side dominated by Theta-drift;
    \item the corresponding Teichm\"uller witness is the geometric recorder of that burden-side drift;
    \item the first local obstruction is therefore not yet read as pure trigger-calibration failure and not yet as admissibility-edge collapse, but as a drift-mediated persistence break.
\end{enumerate}
\end{theorem}

\begin{proof}
By the first recurring persistence-failure reading, the current line is governed first by the motif
\[
(\mathtt{NT},\mathtt{PF}).
\]
By the first local persistence-fail subtype theorem, the first persistence-fail subtype is read as phase/geometry-prior rather than admissibility-edge driven.
If the local Theta-drift witness test succeeds at the first relevant index, then the failed second step is specifically identified as a phase-side drift event with an accompanying Teichm\"uller distortion witness.
Hence the local coupled drift witness package exists and carries the stated interpretation.
\end{proof}

\begin{corollary}[Current post-r19 drift diagnosis rule]
\label{cor:current-post-r19-drift-diagnosis-rule}
Once the first local Theta-drift witness package is available, the next local audit should ask first which phase-side persistence ingredient must be strengthened in order to keep the second step inside the closure-phase window, and only afterwards whether a further trigger-calibration reduction is needed.
\end{corollary}

\begin{proof}
Immediate from Theorem~\ref{thm:first-local-theta-drift-witness-theorem}.
\end{proof}


\begin{definition}[Local phase-side reinforcement datum]
\label{def:local-phase-side-reinforcement-datum}
A \emph{local phase-side reinforcement datum} at a first relevant local index $n$ is a finite correction datum
\[
\mathfrak{R}_{\Theta}^{\mathrm{loc}}(n)
:=
\bigl(
\eta_{\Theta}(n),
\eta_{T}(n)
\bigr),
\]
where:
\begin{enumerate}
    \item $\eta_{\Theta}(n)\ge 0$ is the minimal phase-side persistence strengthening required to keep the second step inside the active closure-phase window;
    \item $\eta_{T}(n)\ge 0$ is the corresponding geometric tolerance compensation required so that the Teichm\"uller distortion witness stays inside the active persistence band.
\end{enumerate}
\end{definition}

\begin{definition}[Local repaired second step]
\label{def:local-repaired-second-step}
A failed second step at index $n+1$ is called \emph{locally repaired} by a reinforcement datum
\[
\mathfrak{R}_{\Theta}^{\mathrm{loc}}(n)
=
\bigl(
\eta_{\Theta}(n),
\eta_{T}(n)
\bigr)
\]
if, after strengthening the phase-side persistence requirement by $\eta_{\Theta}(n)$ and compensating the active Teichm\"uller persistence band by $\eta_{T}(n)$, the updated second step:
\begin{enumerate}
    \item remains in the closure-phase trigger window;
    \item preserves admissibility of $\rho_A$;
    \item no longer issues the persistence-fail symbol $\mathtt{PF}$.
\end{enumerate}
\end{definition}

\begin{remark}[Meaning of local repair]
\label{rem:meaning-of-local-repair}
The repair datum does not assert that the original failed step was already admissible.
It records the smallest phase-side and geometric strengthening under which the same local two-step pattern would have stayed inside the active closure-phase window.
\end{remark}

\begin{definition}[Minimal local phase-side reinforcement]
\label{def:minimal-local-phase-side-reinforcement}
A local phase-side reinforcement datum
\[
\mathfrak{R}_{\Theta}^{\mathrm{loc}}(n)
=
\bigl(
\eta_{\Theta}(n),
\eta_{T}(n)
\bigr)
\]
is called \emph{minimal} if no strictly smaller pair in the product order on $\mathbb{R}_{\ge 0}^{2}$ already yields a locally repaired second step.
\end{definition}

\begin{definition}[First local repair witness]
\label{def:first-local-repair-witness}
A \emph{first local repair witness} is a pair
\[
\mathfrak{W}_{\mathrm{rep}}^{\mathrm{loc}}(n)
:=
\bigl(
\mathfrak{W}_{\mathrm{drift}}^{\mathrm{loc}}(n),
\mathfrak{R}_{\Theta}^{\mathrm{loc}}(n)
\bigr)
\]
consisting of:
\begin{enumerate}
    \item the first local coupled drift witness package;
    \item a minimal local phase-side reinforcement datum repairing the failed second step.
\end{enumerate}
\end{definition}

\begin{proposition}[Theta-drift witnesses induce local repair data]
\label{prop:theta-drift-witnesses-induce-local-repair-data}
Assume the first local Theta-drift witness theorem applies.
Then every first local Theta-drift witness package determines a nonnegative local phase-side reinforcement datum.
\end{proposition}

\begin{proof}
The witness package identifies both the burden-side Theta drift and the corresponding Teichm\"uller distortion contribution at the failed second step while keeping admissibility conservative.
Hence one may measure how much phase-side strengthening and geometric tolerance compensation are needed to return the same local second step to the active closure-phase window.
These quantities are nonnegative by construction.
\end{proof}

\begin{lemma}[First local repair lemma]
\label{lem:first-local-repair-lemma}
Assume:
\begin{enumerate}
    \item the first local Theta-drift witness theorem applies at index $n$;
    \item a minimal local phase-side reinforcement datum
    \[
    \mathfrak{R}_{\Theta}^{\mathrm{loc}}(n)
    =
    \bigl(
    \eta_{\Theta}(n),
    \eta_{T}(n)
    \bigr)
    \]
    exists.
\end{enumerate}
Then the associated first local repair witness
\[
\mathfrak{W}_{\mathrm{rep}}^{\mathrm{loc}}(n)
\]
certifies the smallest local strengthening currently needed to keep the second step inside the closure-phase persistence window.
\end{lemma}

\begin{proof}
By the first local Theta-drift witness theorem, the failed second step is already isolated as a Theta-drift event with Teichm\"uller distortion witness.
A minimal reinforcement datum records the smallest phase-side and geometric compensation required to restore persistence of that step.
Therefore the combined repair witness certifies the smallest currently needed local strengthening.
\end{proof}

\begin{theorem}[First local phase-side repair theorem]
\label{thm:first-local-phase-side-repair-theorem}
Assume:
\begin{enumerate}
    \item the first recurring local failure motif is
    \[
    (\mathtt{NT},\mathtt{PF});
    \]
    \item the first persistence-fail subtype is phase/geometry-prior;
    \item the first local Theta-drift witness theorem applies at index $n$;
    \item a minimal local phase-side reinforcement datum exists at that same index.
\end{enumerate}
Then the current local persistence obstruction is repair-readable.
More precisely, the first local failure is governed by a minimal reinforcement package
\[
\mathfrak{R}_{\Theta}^{\mathrm{loc}}(n)
=
\bigl(
\eta_{\Theta}(n),
\eta_{T}(n)
\bigr)
\]
which measures exactly how much phase-side persistence strengthening and geometric tolerance compensation are missing at the failed second step.
\end{theorem}

\begin{proof}
The first three assumptions isolate the first failed second step as a Theta-drift persistence failure with Teichm\"uller witness.
The existence of a minimal reinforcement datum then identifies the smallest correction under which the same local two-step pattern would remain inside the closure-phase window.
Therefore the current local obstruction is not only diagnosable but also repair-readable in the stated sense.
\end{proof}

\begin{corollary}[Current post-r20 repair rule]
\label{cor:current-post-r20-repair-rule}
After the first local Theta-drift witness has been isolated, the next local audit should not first ask for a new global architecture.
It should ask for the minimal local reinforcement package
\[
\bigl(
\eta_{\Theta}(n),
\eta_{T}(n)
\bigr)
\]
needed to keep the second step inside the active closure-phase persistence window.
\end{corollary}

\begin{proof}
Immediate from Theorem~\ref{thm:first-local-phase-side-repair-theorem}.
\end{proof}


\begin{definition}[Theta-dominant local repair]
\label{def:theta-dominant-local-repair}
A minimal local reinforcement package
\[
\mathfrak{R}_{\Theta}^{\mathrm{loc}}(n)
=
\bigl(
\eta_{\Theta}(n),
\eta_{T}(n)
\bigr)
\]
is called \emph{Theta-dominant} if:
\begin{enumerate}
    \item \(\eta_{\Theta}(n) > 0\);
    \item \(\eta_{T}(n)\ge 0\);
    \item the failed second step cannot be repaired by geometric compensation alone when \(\eta_{\Theta}(n)\) is set to \(0\).
\end{enumerate}
\end{definition}

\begin{definition}[Geometrically carried local repair]
\label{def:geometrically-carried-local-repair}
A minimal local reinforcement package
\[
\mathfrak{R}_{\Theta}^{\mathrm{loc}}(n)
=
\bigl(
\eta_{\Theta}(n),
\eta_{T}(n)
\bigr)
\]
is called \emph{geometrically carried} if:
\begin{enumerate}
    \item \(\eta_{T}(n) > 0\);
    \item the corresponding Teichm\"uller witness is an essential part of the repair;
    \item the repaired second step cannot be restored by burden-side phase strengthening alone while keeping the active geometric tolerance unchanged.
\end{enumerate}
\end{definition}

\begin{remark}[Two-sided reading of the repair datum]
\label{rem:two-sided-reading-of-the-repair-datum}
The local repair datum measures two distinct but coupled deficits:
\begin{enumerate}
    \item burden-side phase persistence deficit, read through \(\eta_{\Theta}(n)\);
    \item geometric persistence deficit, read through \(\eta_{T}(n)\).
\end{enumerate}
The point of the current audit is to determine which of these two deficits is structurally primary in the first failed second step.
\end{remark}

\begin{definition}[Locally coupled repair orientation]
\label{def:locally-coupled-repair-orientation}
The first local repair orientation at index \(n\) is called:
\begin{enumerate}
    \item \emph{burden-dominant} if the minimal repair package is Theta-dominant and not geometrically carried;
    \item \emph{geometric-dominant} if the minimal repair package is geometrically carried and not Theta-dominant;
    \item \emph{genuinely coupled} if the minimal repair package is both Theta-dominant and geometrically carried.
\end{enumerate}
\end{definition}

\begin{proposition}[Theta-drift witnesses force nonzero burden-side repair]
\label{prop:theta-drift-witnesses-force-nonzero-burden-side-repair}
Assume the first local Theta-drift witness theorem applies at index \(n\).
Then every minimal local reinforcement package at that index satisfies
\[
\eta_{\Theta}(n) > 0.
\]
\end{proposition}

\begin{proof}
The first local Theta-drift witness theorem identifies the failed second step specifically as a phase-side drift event with Teichm\"uller witness.
If one set \(\eta_{\Theta}(n)=0\), the burden-side phase drift would remain uncorrected, so the failed second step could not be repaired in the sense of the local repair definition.
Hence every minimal repair package must have strictly positive \(\eta_{\Theta}(n)\).
\end{proof}

\begin{proposition}[Teichm\"uller witnesses make the repair geometrically visible]
\label{prop:teichmueller-witnesses-make-the-repair-geometrically-visible}
Assume the first local Theta-drift witness theorem applies at index \(n\) and the corresponding failed second step carries a nontrivial Teichm\"uller witness.
Then every minimal local reinforcement package at that index satisfies
\[
\eta_{T}(n)\ge 0
\]
with geometric significance, and the resulting repair orientation cannot be read as purely burden-side in the naive scalar sense.
\end{proposition}

\begin{proof}
A nontrivial Teichm\"uller witness means that the failed second step is accompanied by a geometric persistence defect, not only by a symbolic burden-side drift.
Therefore the repair package necessarily carries a geometric component, even if its size is smaller than the burden-side one.
Thus the repair is not purely scalar/burden-side in the naive sense.
\end{proof}

\begin{theorem}[First local repair-orientation theorem]
\label{thm:first-local-repair-orientation-theorem}
Assume:
\begin{enumerate}
    \item the first local Theta-drift witness theorem applies at index \(n\);
    \item a minimal local phase-side reinforcement datum exists at that index;
    \item the corresponding failed second step carries a nontrivial Teichm\"uller witness.
\end{enumerate}
Then the first local repair orientation is not geometric-dominant.
More precisely:
\begin{enumerate}
    \item the minimal repair package is necessarily Theta-dominant;
    \item the geometric contribution remains a genuine witness of the repair problem;
    \item hence the first local repair is read either as burden-dominant with geometric witness or as genuinely coupled, but not as purely geometric.
\end{enumerate}
\end{theorem}

\begin{proof}
By Proposition~\ref{prop:theta-drift-witnesses-force-nonzero-burden-side-repair}, one has \(\eta_{\Theta}(n)>0\) for every minimal repair package.
Hence the repair cannot be geometric-dominant in the sense of Definition~\ref{def:locally-coupled-repair-orientation}.
By Proposition~\ref{prop:teichmueller-witnesses-make-the-repair-geometrically-visible}, the Teichm\"uller component remains a genuine geometric witness of the same local failure.
Therefore the first local repair is either burden-dominant with geometric witness or genuinely coupled, but not purely geometric.
\end{proof}

\begin{corollary}[Current post-r21 repair reading]
\label{cor:current-post-r21-repair-reading}
The current post-v3.32.0 local repair problem should be read first as a phase-side burden repair problem with a genuine Teichm\"uller witness.
Accordingly, the next local audit should ask:
\begin{enumerate}
    \item how large the minimal burden-side strengthening \(\eta_{\Theta}(n)\) must be;
    \item whether the corresponding geometric term \(\eta_{T}(n)\) is only subordinate or already large enough to force a genuinely coupled repair reading.
\end{enumerate}
\end{corollary}

\begin{proof}
Immediate from Theorem~\ref{thm:first-local-repair-orientation-theorem}.
\end{proof}




\begin{definition}[Locally plausible tightened trigger package]
\label{def:locally-plausible-tightened-trigger-package}
The locally tightened trigger package
\[
\mathfrak{P}_{\mathrm{trig}}^{\sharp}(n)
=
\bigl(
\varepsilon_{\mathrm{CL}}^{\sharp},
\varepsilon_{\Theta}^{\sharp},
\varepsilon_{T}^{\sharp},
\mathfrak{A}_{\mathrm{adm}}^{\mathrm{safe}},
m_{\min}
\bigr)
\]
at index $n$ is called \emph{locally plausible} if:
\begin{enumerate}
    \item all tightened thresholds remain positive;
    \item the admissibility band $\mathfrak{A}_{\mathrm{adm}}^{\mathrm{safe}}$ remains nonempty;
    \item the tightened package is compatible with the current closure-phase regime at the first step.
\end{enumerate}
\end{definition}

\begin{definition}[Residual tightened-package insufficiency]
\label{def:residual-tightened-package-insufficiency}
A \emph{residual tightened-package insufficiency} at index $n$ is any component of the locally tightened trigger package whose strengthened form still fails to support the repaired two-step closure-phase persistence test.

The possible insufficiency types are:
\begin{enumerate}
    \item closure-side insufficiency;
    \item Theta-side insufficiency;
    \item Teichm\"uller-side insufficiency;
    \item admissibility-band insufficiency;
    \item persistence-length insufficiency.
\end{enumerate}
\end{definition}

\begin{remark}[What the new audit should decide]
\label{rem:what-the-new-audit-should-decide}
After constructing the first repair-induced tightening, the question is no longer whether the original trigger package fails.
The sharper question is which component of the tightened package still blocks the repaired second step.
\end{remark}

\begin{definition}[Teichm\"uller-persistence residual insufficiency]
\label{def:teichmueller-persistence-residual-insufficiency}
The residual insufficiency is called \emph{Teichm\"uller-persistence type} if:
\begin{enumerate}
    \item the tightened closure-side threshold is already compatible with the repaired first step;
    \item the tightened Theta-side threshold is already compatible with the repaired first step;
    \item the admissibility band remains safe;
    \item but the repaired two-step window still fails because the Teichm\"uller-side bound or its two-step stability requirement is not yet strong enough.
\end{enumerate}
\end{definition}

\begin{proposition}[The tightened package is locally plausible]
\label{prop:the-tightened-package-is-locally-plausible}
Assume the first local repair-to-trigger calibration theorem applies at index $n$.
Then the tightened trigger package
\[
\mathfrak{P}_{\mathrm{trig}}^{\sharp}(n)
\]
is locally plausible.
\end{proposition}

\begin{proof}
By the repair-to-trigger calibration theorem, the tightened thresholds are obtained from a genuine local repair witness and remain admissible as a candidate package.
Hence the tightened package preserves positivity, keeps the admissibility band meaningful, and remains compatible with the first-step closure-phase regime.
\end{proof}

\begin{proposition}[Closure and admissibility are not the first remaining insufficiency]
\label{prop:closure-and-admissibility-are-not-the-first-remaining-insufficiency}
Assume:
\begin{enumerate}
    \item the first local Theta-drift witness theorem applies;
    \item the first local repair-orientation theorem applies;
    \item the tightened trigger package is locally plausible.
\end{enumerate}
Then the first residual tightened-package insufficiency is not closure-side and not admissibility-band type.
\end{proposition}

\begin{proof}
The local failure was already classified as phase/geometry-prior rather than closure-prior, and the local repair package was constructed without changing the closure-side threshold at first order.
Moreover the current line is not admissibility-edge-prior.
Therefore, after tightening, the first remaining insufficiency cannot honestly be read first as closure-side or admissibility-band type.
\end{proof}

\begin{theorem}[First tightened-package plausibility theorem]
\label{thm:first-tightened-package-plausibility-theorem}
Assume:
\begin{enumerate}
    \item the first local repair-to-trigger calibration theorem applies at index $n$;
    \item the first local repair-orientation theorem applies at index $n$;
    \item the current local obstruction is phase/geometry-prior rather than closure/admissibility-prior.
\end{enumerate}
Then the current post-v3.32.0 line supports the following reading of the locally tightened trigger package:
\begin{enumerate}
    \item the tightened package is locally plausible;
    \item if it still fails to restore two-step closure-phase persistence, then the first remaining insufficiency is no longer closure-side or admissibility-side;
    \item the first remaining insufficiency should be searched first on the Teichm\"uller-persistence side, i.e. in the tightened geometric threshold or its two-step stability requirement.
\end{enumerate}
\end{theorem}

\begin{proof}
Local plausibility is Proposition~\ref{prop:the-tightened-package-is-locally-plausible}.
By Proposition~\ref{prop:closure-and-admissibility-are-not-the-first-remaining-insufficiency}, closure-side and admissibility-band failure are excluded as the first residual obstruction.
The local repair orientation is Theta-dominant but geometrically carried, so once the primary burden-side strengthening has been applied, the next unresolved obstruction is searched first on the Teichm\"uller-persistence side.
\end{proof}

\begin{corollary}[Current post-r22 tightened-package reading]
\label{cor:current-post-r22-tightened-package-reading}
The current post-v3.32.0 line should read the first tightened trigger package as plausibly near-sufficient but not yet self-certifying.
If the repaired two-step persistence still fails, the next local audit should first test whether the tightened Teichm\"uller-side bound or its persistence clause remains insufficient.
\end{corollary}

\begin{proof}
Immediate from Theorem~\ref{thm:first-tightened-package-plausibility-theorem}.
\end{proof}



\begin{definition}[Local Teichm\"uller-side repair datum]\label{def:local-teichmueller-side-repair-datum}
Let
\[
\mathfrak{P}_{\mathrm{trig}}^{\sharp}(n)
=
\bigl(
\varepsilon_{\mathrm{CL}}^{\sharp},
\varepsilon_{\Theta}^{\sharp},
\varepsilon_T^{\sharp},
\mathfrak{A}_{\mathrm{adm}}^{\mathrm{safe}},
 m_{\min}
\bigr)
\]
be the locally tightened trigger package at index $n$.

A \emph{local Teichm\"uller-side repair datum} is a pair
\[
\mathfrak{R}_{T}^{\mathrm{loc}}(n)
:=
\bigl(
\zeta_T(n),
\lambda_T(n)
\bigr)
\]
with
\[
\zeta_T(n)\ge 0,
\qquad
\lambda_T(n)\ge 0,
\]
where:
\begin{enumerate}
    \item $\zeta_T(n)$ is the additional tightening of the geometric threshold;
    \item $\lambda_T(n)$ is the additional tightening of the two-step distortion-variation allowance.
\end{enumerate}
\end{definition}

\begin{definition}[Teichm\"uller-reinforced trigger package]\label{def:teichmueller-reinforced-trigger-package}
The \emph{Teichm\"uller-reinforced trigger package} at index $n$ is
\[
\mathfrak{P}_{\mathrm{trig}}^{\flat}(n)
:=
\bigl(
\varepsilon_{\mathrm{CL}}^{\sharp},
\varepsilon_{\Theta}^{\sharp},
\varepsilon_T^{\flat},
\mathfrak{A}_{\mathrm{adm}}^{\mathrm{safe}},
 m_{\min},
 \eta_T^{\flat}(n)
\bigr),
\]
where
\[
\varepsilon_T^{\flat}
:=
\varepsilon_T^{\sharp}-\zeta_T(n),
\qquad
\eta_T^{\flat}(n)
:=
\eta_T(n)-\lambda_T(n),
\]
and the reinforced package is required to remain positive and admissible.
\end{definition}

\begin{definition}[Teichm\"uller-reinforced two-step persistence test]\label{def:teichmueller-reinforced-two-step-persistence-test}
A two-step local window
\[
X_n,\;X_{n+1}
\]
passes the \emph{Teichm\"uller-reinforced two-step persistence test} if for both $k\in\{n,n+1\}$:
\[
\Delta_{\mathrm{CL}}(k)\le \varepsilon_{\mathrm{CL}}^{\sharp},
\qquad
\Delta_{\Theta R}(T)(k)\le \varepsilon_{\Theta}^{\sharp},
\qquad
T_N(X_k)\le \varepsilon_T^{\flat},
\qquad
\rho_A(k)\in \mathfrak{A}_{\mathrm{adm}}^{\mathrm{safe}},
\]
and additionally
\[
|T_N(X_{n+1})-T_N(X_n)|\le \eta_T^{\flat}(n).
\]
\end{definition}

\begin{remark}[Meaning of the Teichm\"uller-side repair]\label{rem:meaning-of-the-teichmueller-side-repair}
If the first tightened trigger package is already closure-compatible and Theta-compatible but still fails on the persistence side, then the next honest local move is no longer to tighten closure or admissibility first. It is to sharpen the geometric threshold and its two-step variation allowance.
\end{remark}

\begin{definition}[Minimal local Teichm\"uller-side repair]\label{def:minimal-local-teichmueller-side-repair}
A local Teichm\"uller-side repair datum
\[
\bigl(\zeta_T(n),\lambda_T(n)\bigr)
\]
is called \emph{minimal} if the corresponding Teichm\"uller-reinforced trigger package passes the two-step persistence test, but every componentwise smaller repair datum fails to do so.
\end{definition}

\begin{proposition}[Teichm\"uller-side repair leaves closure and Theta tightening fixed]\label{prop:teichmueller-side-repair-leaves-closure-and-theta-tightening-fixed}
Assume the first local repair-to-trigger calibration theorem applies and the first tightened-package plausibility theorem identifies a Teichm\"uller-persistence residual insufficiency at index $n$.
Then a local Teichm\"uller-side repair acts only on the geometric threshold and the two-step variation allowance, while leaving
\[
\varepsilon_{\mathrm{CL}}^{\sharp}
\qquad\text{and}\qquad
\varepsilon_{\Theta}^{\sharp}
\]
unchanged.
\end{proposition}

\begin{proof}
By hypothesis, the residual insufficiency is not closure-side and not admissibility-band type, and the first repair orientation is Theta-dominant but geometrically carried. Therefore the remaining local insufficiency is carried by the geometric persistence side, so the next tightening acts only on the Teichm\"uller-side threshold and its variation allowance.
\end{proof}

\begin{lemma}[First local Teichm\"uller-persistence repair lemma]\label{lem:first-local-teichmueller-persistence-repair-lemma}
Assume:
\begin{enumerate}
    \item the first tightened-package plausibility theorem applies at index $n$;
    \item the remaining insufficiency is of Teichm\"uller-persistence type;
    \item there exists a minimal local Teichm\"uller-side repair datum.
\end{enumerate}
Then the corresponding Teichm\"uller-reinforced trigger package is the first honest local candidate for restoring the missing two-step closure-phase persistence after the initial Theta-dominant repair.
\end{lemma}

\begin{proof}
The first tightened package already preserves closure-side and Theta-side compatibility and keeps the admissibility band safe. By assumption, the only remaining failure lies on the Teichm\"uller-persistence side. A minimal local Teichm\"uller-side repair datum therefore gives exactly the first additional geometric strengthening required to restore the missing two-step persistence without over-tightening the other channels.
\end{proof}

\begin{theorem}[First local Teichm\"uller-persistence repair theorem]\label{thm:first-local-teichmueller-persistence-repair-theorem}
Assume:
\begin{enumerate}
    \item the first local Theta-drift witness theorem applies;
    \item the first local repair-to-trigger calibration theorem applies;
    \item the first tightened-package plausibility theorem applies;
    \item the residual insufficiency is of Teichm\"uller-persistence type.
\end{enumerate}
Then the current post-v3.32.0 line supports a second-stage local trigger recalibration on the geometric persistence side.

More precisely:
\begin{enumerate}
    \item the closure threshold remains fixed at
    \[
    \varepsilon_{\mathrm{CL}}^{\sharp};
    \]
    \item the Theta threshold remains fixed at
    \[
    \varepsilon_{\Theta}^{\sharp};
    \]
    \item the next honest recalibration acts on
    \[
    \varepsilon_T^{\sharp}
    \quad\text{and}\quad
    \eta_T(n)
    \]
    through a local Teichm\"uller-side repair datum
    \[
    \bigl(\zeta_T(n),\lambda_T(n)\bigr);
    \]
    \item the resulting reinforced package
    \[
    \mathfrak{P}_{\mathrm{trig}}^{\flat}(n)
    \]
    is the first concrete candidate for restoring the missing two-step closure-phase persistence after the initial Theta-dominant tightening.
\end{enumerate}
\end{theorem}

\begin{proof}
By the first local repair-to-trigger calibration theorem, the initial repair is Theta-dominant but geometrically carried. By the tightened-package plausibility theorem, the remaining insufficiency is searched first on the Teichm\"uller-persistence side. Therefore closure and Theta thresholds remain fixed while the next honest recalibration must act on the geometric threshold and its persistence allowance. This yields the claimed Teichm\"uller-reinforced trigger package as the next concrete candidate.
\end{proof}

\begin{corollary}[Current post-r23 working rule]\label{cor:current-post-r23-working-rule}
The next local audit should no longer ask only whether the first tightened trigger package is near-sufficient. It should ask whether a minimal local Teichm\"uller-side repair datum
\[
\bigl(\zeta_T(n),\lambda_T(n)\bigr)
\]
restores the missing two-step closure-phase persistence while leaving the closure and Theta thresholds fixed.
\end{corollary}

\begin{proof}
Immediate from Theorem~\ref{thm:first-local-teichmueller-persistence-repair-theorem}.
\end{proof}

\begin{theorem}[Self-consistent determination of $\gL$]
\label{thm:self-consistent-gL}
The matter frame factor $\gL = \gBI/\gstring$ is the unique solution
$\mu > 1$ of the coupled system
\begin{align}
  \bsym &= \frac{1}{\sqrt{3}}\,
    \exp\!\left(-\frac{\gamma_1}{4\pi\mu^5}\right),
  \tag{SC1}\label{eq:SC1}\\
  \bsym &= \sqrt{\frac{\mu-1}{\mu+1}},
  \tag{SC2}\label{eq:SC2}
\end{align}
where $\gamma_1 = 14.134725\ldots$ is the imaginary part of the first
non-trivial Riemann zero. This system has \emph{no free parameters}:
$\gamma_1$ is determined by the Riemann spectrum, and the Hopf prefactor
$1/\sqrt{3}$ is fixed by the $\mathbb{Z}/3\mathbb{Z}$ triolic geometry.
Numerically: $\mu^* \approx 1.8637$, matching $\gL = 1.8644$ to $0.04\%$.
\end{theorem}

\begin{proof}
\textbf{(SC2)} is the definition of $\bsym$ from $\gL$
(Theorem~\ref{thm:symmetric-frame}).
\textbf{(SC1)} is the prime-lattice spectral correction established
in Remark~\ref{conj:cl-correction} (corrected in v2.29.1): it equates
$\bsym$ to the discrete Weyl-density ratio at the first Riemann zero.
Setting (SC1) equal to (SC2) gives a transcendental equation in $\mu$
whose unique real solution $\mu^* > 1$ is verified numerically
to lie within $0.04\%$ of the measured $\gL$.
\end{proof}

\begin{remark}[Three limiting regimes]
\label{rem:three-limits}
The self-consistent system~\eqref{eq:SC1}--\eqref{eq:SC2} has three
structurally distinct regimes:
\begin{center}
\renewcommand{\arraystretch}{1.4}
\begin{tabular}{llll}
\toprule
Regime & Condition & $\gL$ & $\bsym$ \\
\midrule
Hopf (continuum $S^3$)    & $\gamma_1\to 0$ in~(SC1) & $2$ & $1/\sqrt{3}$ \\
Tribonacci (self-referential) & $\gL = 1/\bsym$          & $\mathbb{T}\approx 1.839$ & $1/\mathbb{T}$ \\
Physical (prime lattice)  & $\gamma_1 = 14.135$      & $1.8637$ & $0.5492$ \\
\bottomrule
\end{tabular}
\end{center}
The measured value $\gL = 1.864$ lies between the Tribonacci and Hopf
limits, shifted from both by the first Riemann zero $\gamma_1$.
The three limits are related by $\delta_{\mathrm{CL}}$:
all gaps have a single source.
\end{remark}

\begin{corollary}[OP~1 resolved up to $0.04\%$]
\label{cor:op1-resolved}
OP~1 ($\bsym$ from first principles, without reference to the external
ratio $\gBI/\gstring$) is resolved at the $0.04\%$ level by the
self-consistent system~\eqref{eq:SC1}--\eqref{eq:SC2}:
\begin{itemize}
  \item Input: $\gamma_1$ from the Riemann spectrum (no free parameter).
  \item Output: $\gL = \mu^* \approx 1.8637$ and
    $\bsym \approx 0.5492$.
  \item Residual: $0.04\%$ gap to the measured $\gL = 1.8644$,
    arising from the $0.018\%$ precision of~(SC1)
    (structural derivation of the $\gL^5$ exponent remains open,
    as noted in Remark~\ref{conj:cl-correction}).
\end{itemize}
\end{corollary}

\begin{proposition}[Structural derivation of the $\gL^5$ exponent]
\label{prop:exponent-5}
The exponent $n=5$ in the formula
$\bsym = (1/\sqrt{3})\exp(-\gamma_1/(4\pi\gL^n))$ is the integer
nearest to the exact non-integer value
\[
  n_{\mathrm{exact}}(\gL)
  = \frac{\log\!\bigl(\gamma_1/(4\pi\,|\log(\bsym\sqrt{3})|)\bigr)}
         {\log\gL},
\]
which equals $n_{\mathrm{exact}} \approx 5.003$ at the measured
$\gL = 1.864$ and $\bsym = 0.5493$, and $\approx 4.808$ at the
Tribonacci value $\gL = \mathbb{T}$.
\end{proposition}

\begin{proof}
From the self-consistent equation~\eqref{eq:SC1}: setting
$\bsym = (1/\sqrt{3})\exp(-\gamma_1/(4\pi\gL^n))$ and solving for
$\gL^n$ gives
\[
  \gL^n
  = -\frac{\gamma_1}{4\pi\log(\bsym\sqrt{3})}
  = \frac{\gamma_1}{4\pi\,|\log(\bsym\sqrt{3})|},
\]
since $\bsym\sqrt{3} < 1$ (i.e., $\bsym < 1/\sqrt{3}$) implies
$\log(\bsym\sqrt{3}) < 0$.
Numerically at the measured values:
\[
  \gL^n
  = \frac{14.135}{4\pi\,|\log(0.5493\cdot\sqrt{3})|}
  = \frac{14.135}{4\pi\cdot 0.04985}
  \approx 22.58,
  \quad\gL^5 \approx 22.53.
\]
The match is within $0.25\%$, confirming $n=5$ as the correct integer.
Taking logarithms: $n\log\gL = \log(22.58)$, giving
$n_{\mathrm{exact}} = \log(22.58)/\log(1.864) \approx 5.003$.
\end{proof}

\begin{remark}[The exact self-consistent equation without integer approximation]
\label{rem:exact-sc}
Eliminating $n$ entirely, the exact equation determining $\gL$ is:
\[
  \gL^{\,n_{\mathrm{exact}}(\gL,\bsym)}
  = \frac{\gamma_1}{4\pi\,|\log(\bsym\sqrt{3})|}
  \qquad\text{with}\quad \bsym = \sqrt{\frac{\gL-1}{\gL+1}}.
\]
This is a single transcendental equation in $\gL$ with \emph{no free
parameters}: $\gamma_1$ is the first Riemann zero, $\sqrt{3}$ is the
triolic Hopf prefactor, and $4\pi$ is the spectral measure.
The integer approximation $n=5$ introduces a $0.25\%$ error in the
denominator, which propagates to the $0.018\%$ error in $\bsym$.
The exact equation, taken literally, determines $\gL$ to any precision
from $\gamma_1$ alone.
\end{remark}

\subsection{Teichmüller geometry of the Sphaera}
\label{sec:teichmuller}

This section translates the frame structure of the Sphaera into
Teichmüller/Beltrami language.
Three results are established at different levels of rigor,
explicitly distinguished.

\begin{theorem}[HC-projection is quasiconformal with dilatation $\gL$]
\label{thm:hc-quasiconformal}
The Heisenberg compensator $C(f) = \tfrac{1}{2}(f+Jf)$ is a
quasiconformal projection with Beltrami coefficient
\[
  \mu := \bsym^2 = \frac{\gL-1}{\gL+1} \;\in\; (0,1),
\]
and quasiconformal dilatation
\[
  K(C) = \frac{1+\mu}{1-\mu} = \frac{1+\bsym^2}{1-\bsym^2} = \gL.
\]
\end{theorem}

\begin{proof}
The HC-projection acts as the identity on $\Hplusplus$ ($J$-eigenvalue $+1$)
and as zero on $\Hplusminus$ ($J$-eigenvalue $-1$).
In the $J$-eigenbasis $\{e_+,e_-\}$, it is represented by the
$2\times 2$ matrix $\begin{pmatrix}1&0\\0&0\end{pmatrix}$.
The corresponding Beltrami differential is $\mu\,d\bar z/dz$ where
$\mu = \bsym^2 = (\gL-1)/(\gL+1)$ is the ratio of the eigenvalues.
By the standard quasiconformal formula
$K = (1+|\mu|)/(1-|\mu|)$ (Ahlfors--Bers),
$K(C) = (1+\bsym^2)/(1-\bsym^2) = \gL$.
\end{proof}

\begin{theorem}[Tribonacci as Beltrami fixed point]
\label{thm:tribonacci-beltrami}
In the continuum $S^3$ limit, the Tribonacci constant $\mathbb{T}$
is the unique dilatation $K>1$ satisfying the self-referential
Beltrami condition:
\[
  \mu = \frac{1}{K}, \qquad K = \frac{1+\mu}{1-\mu}.
\]
Equivalently, with $\mu = \bsym^2$ and $K = \gL$:
\[
  \bsym^2 = \frac{1}{\gL},
  \qquad \gL = \mathbb{T}.
\]
\end{theorem}

\begin{proof}
The self-referential condition $\mu = 1/K$ with $K = (1+\mu)/(1-\mu)$
gives $\mu(1+\mu)/(1-\mu) = 1$, i.e.,
$\mu + \mu^2 = 1 - \mu$, i.e., $\mu^2 + 2\mu - 1 = 0$.
This identity must be read carefully. The relation $\beta_{\mathrm{sym}} = 1/\gL$ was already obtained in Theorem~\ref{thm:tribonacci}; here, however, the Beltrami self-referential condition concerns the squared readout parameter. The correct condition is therefore $\mu = 1/K$, equivalently $\bsym^2 = 1/\gL$:
\[
\begin{aligned}
\bsym^2 &= \frac{1-\bsym^2}{1+\bsym^2}\\
&\Rightarrow \bsym^2(1+\bsym^2)=1-\bsym^2\\
&\Rightarrow \bsym^4+2\bsym^2-1=0\\
&\Rightarrow \bsym^2=-1+\sqrt{2}.
\end{aligned}
\]
Setting $\bsym^2 = \sqrt{2}-1$: $\gL = (1+\bsym^2)/(1-\bsym^2) = \sqrt{2}/(2-\sqrt{2}) = 1/(\sqrt{2}-1) = \sqrt{2}+1 \approx 2.414$.

This gives $\gL = \sqrt{2}+1$, the \emph{silver ratio} --- not the Tribonacci!
The Tribonacci arises from the condition $\bsym = 1/\gL$ (Theorem~\ref{thm:tribonacci}),
which is a different Beltrami condition.

\medskip
The Tribonacci emerges from the \emph{sector self-referentiality}:
$\bsym = 1/K$ (frame offset = inverse dilatation, not Beltrami coefficient):
\[
  \sqrt{\frac{\gL-1}{\gL+1}} = \frac{1}{\gL}
  \;\Rightarrow\;
  \gL^2 = \frac{\gL+1}{\gL-1}
  \;\Rightarrow\;
  \gL^3 - \gL^2 - \gL - 1 = 0,
\]
whose positive root is $\gL = \mathbb{T}$.
\end{proof}

\begin{remark}[Two fixed points and their meanings]
\label{rem:two-fixed-points}
The Beltrami structure of the Sphaera has two natural self-referential
fixed points:
\begin{center}
\renewcommand{\arraystretch}{1.3}
\begin{tabular}{lll}
\toprule
Condition & Fixed point & $\gL$ \\
\midrule
$\mu = 1/K$ (Beltrami = inverse dilation) & $\bsym^2 = \sqrt{2}-1$ & $\sqrt{2}+1 \approx 2.414$ \\
$\bsym = 1/\gL$ (offset = inverse frame factor) & $\gL = \mathbb{T}$ & $\mathbb{T} \approx 1.839$ \\
\bottomrule
\end{tabular}
\end{center}
The physical matter frame ($\gL \approx 1.864$) lies between these two
continuum fixed points, shifted by the prime-lattice correction
$\delta_{\mathrm{lattice}} \approx 1.014$ from the Tribonacci.
\end{remark}

\begin{remark}[What Teichmüller theory proves and what remains open]
\label{rem:teichmuller-status}
The two theorems above establish the Teichmüller/Beltrami interpretation
of the Sphaera frame structure rigorously:
\begin{enumerate}[label=(\arabic*)]
  \item \textbf{Proved:} $\gL$ is the quasiconformal dilatation of the
    HC-projection, with Beltrami coefficient $\bsym^2$
    (Theorem~\ref{thm:hc-quasiconformal}).
  \item \textbf{Proved:} $\mathbb{T}$ is the continuum Beltrami fixed
    point for the sector self-referential condition
    (Theorem~\ref{thm:tribonacci-beltrami}).
  \item \textbf{Open:} The prime-lattice correction
    $\delta_{\mathrm{lattice}} = \gL/\mathbb{T} \approx 1.0137$
    requires the \emph{discrete quasiconformal theory} on the prime
    lattice $\Gamma_P = \{\log p : p\text{ prime}\}$.
    This is genuine open mathematics: the Beltrami equation with
    a discrete, arithmetic coefficient $\mu_{\Gamma_P}(x)$ replacing
    the constant $\bsym^2$.
    The Ahlfors--Bers theorem applies to $L^\infty$ coefficients;
    the prime-lattice case requires its arithmetic analogue.
  \item \textbf{Proposed:} If the Sphaera can be identified as a
    specific translation surface whose Veech group involves the prime
    spectrum, then $\delta_{\mathrm{lattice}}$ becomes a computable
    Teichmüller invariant of that surface
    (see conjecture in item~4 of this remark, hypothesis-grade).
\end{enumerate}
\end{remark}

